\documentclass[11pt,a4paper,spanish]{book}
\usepackage{estilo_unir-1}
\usepackage{apacite}
\usepackage{hyperref}

\numberwithin{equation}{chapter}
\numberwithin{figure}{chapter}

%---------------------------
% Datos del trabajo
%---------------------------
\title{EVMAudit: Librería multiherramienta para la detección automatizada de vulnerabilidades en contratos inteligentes de Ethereum}
\titulacion{Máster Universitario en Ciberseguridad}
\author{Adrián Moreno Martín, Daniel Rovira Martínez,Paula Suárez Prieto }
\date{Curso 2025-2026}
\director{Friman Sánchez Castaño}
\nombreciudad{Madrid}

\begin{document}

\pagestyle{fancy}

\renewcommand{\listfigurename}{Índice de Ilustraciones}
\renewcommand{\listtablename}{Índice de Tablas}
\renewcommand{\contentsname}{Índice de Contenidos}
\renewcommand{\figurename}{Figura}
\renewcommand{\tablename}{Tabla}

\maketitle

\frontmatter


\input{capitulos/0.1. Abstract}


\section{Organización del trabajo en
grupo}\label{organizaciuxf3n-del-trabajo-en-grupo}

\subsection{Partes que aborda el TFE, distribución y estructura de la
memoria}\label{partes-que-aborda-el-tfe-distribuciuxf3n-y-estructura-de-la-memoria}

Organización del trabajo en grupo - Desarrollo de la memoria

{\def\LTcaptype{none} % do not increment counter
\begin{longtable}[]{@{}
  >{\raggedright\arraybackslash}p{(\linewidth - 2\tabcolsep) * \real{0.5000}}
  >{\raggedright\arraybackslash}p{(\linewidth - 2\tabcolsep) * \real{0.5000}}@{}}
\toprule\noalign{}
\begin{minipage}[b]{\linewidth}\raggedright
Apartado de la memoria
\end{minipage} & \begin{minipage}[b]{\linewidth}\raggedright
Responsables
\end{minipage} \\
\midrule\noalign{}
\endhead
\bottomrule\noalign{}
\endlastfoot
Introducción & Alumno 1, Alumno 2 y Alumno 3 \\
Estado del arte & Alumno 1, Alumno 2 y Alumno 3 \\
Objetivos y metodología de trabajo & Alumno 1, Alumno 2 y Alumno 3 \\
Desarrollo específico de la contribución: parte individual 1 & Alumno
1 \\
Desarrollo específico de la contribución: parte individual 2 & Alumno
2 \\
Desarrollo específico de la contribución: parte individual 3 & Alumno
3 \\
Desarrollo específico de la contribución: parte 4 & Alumno 1, Alumno 2 y
Alumno 3 \\
Conclusiones y trabajo futuro & Alumno 1, Alumno 2 y Alumno 3 \\
\end{longtable}
}

\subsection{Mecanismos de coordinación
empleados}\label{mecanismos-de-coordinaciuxf3n-empleados}

Con el objetivo de garantizar un seguimiento continuo del trabajo y
mantener la coordinación entre los miembros del grupo, se estableció una
reunión interna semanal en la que cada integrante exponía las tareas
realizadas durante el periodo correspondiente y las posibles
dificultades encontradas. Estas reuniones permitían, además, debatir las
siguientes líneas de trabajo y planificar las tareas a desarrollar en
las semanas posteriores.

Como herramientas de comunicación y colaboración se utilizaron
principalmente WhatsApp para la comunicación diaria, Microsoft Teams
para la realización de reuniones telemáticas, un repositorio compartido
en GitHub para la gestión y sincronización del código fuente, y un
documento compartido de Microsoft Word para la elaboración conjunta de
la memoria.

Adicionalmente, tras cada una de las entregas parciales previstas en la
planificación del Trabajo Fin de Máster, se mantuvo una reunión de
seguimiento con el director del trabajo. Estas sesiones permitieron
revisar el progreso realizado, recibir retroalimentación sobre los
resultados obtenidos y definir las acciones necesarias para las
siguientes fases del proyecto.



\newpage
\tableofcontents
\listoffigures
\listoftables


\mainmatter

\chapter{Introducción}
\input{capitulos/1. Introducción}

\chapter{Estado del arte}
\input{capitulos/2. Estado del arte/2.00.Introduccion}
\subsection{2.1. Fundamentos blockchain y contratos
inteligentes}\label{fundamentos-blockchain-y-contratos-inteligentes}

La tecnología~\textbf{\emph{blockchain}}~se define como un sistema de
registro distribuido (\emph{distributed ledger}) que permite almacenar y
gestionar información de forma descentralizada, segura y resistente a
manipulaciones, sin depender de una autoridad central de confianza. Este
paradigma fue introducido por~Satoshi Nakamoto~en 2008, en el contexto
del diseño de Bitcoin, donde se propone un sistema de dinero
electrónico~\emph{peer-to-peer}~basado en un libro mayor compartido
entre múltiples nodos de la red. {[}1{]}

En una blockchain, las transacciones se agrupan en estructuras
denominadas \textbf{bloques}, que se enlazan secuencialmente mediante
\textbf{funciones hash criptográficas}. Cada bloque contiene, entre
otros elementos, un conjunto de transacciones validadas, una marca
temporal y el hash del bloque anterior. Este encadenamiento garantiza la
integridad del sistema, ya que cualquier modificación en un bloque
previamente confirmado alteraría su hash, rompiendo la continuidad de la
cadena y permitiendo detectar manipulaciones de forma inmediata. {[}2{]}

Figura 1. Cadena genérica de bloques

Fuente: Blockchain Technology Overview, NISTIR 8202.

Adicionalmente, blockchain emplea criptografía asimétrica para asegurar
la autenticidad e integridad de las transacciones. Cada participante
dispone de un par de claves criptográficas (clave pública y clave
privada), donde la clave privada se utiliza para firmar digitalmente las
transacciones, y la clave pública permite a otros nodos verificar su
validez. Este mecanismo elimina la necesidad de intermediarios de
confianza, trasladando la verificación al propio sistema distribuido.
{[}2{]}

\subsubsection{2.1.1. ~Tipos de redes
blockchain}\label{tipos-de-redes-blockchain}

Las redes blockchain pueden clasificarse en función de su modelo de
acceso y gobernanza, distinguiéndose principalmente entre redes públicas
(\emph{permissionless}) y redes privadas (\emph{permissioned}).

Las~\textbf{\emph{blockchains}} \emph{\textbf{permisionless}}~permiten
la participación abierta de cualquier usuario, que puede leer el estado
del \emph{ledger} y, dependiendo del protocolo de consenso, participar
en la validación de bloques. Este modelo prioriza la descentralización,
la transparencia y la resistencia a la censura. Ejemplos representativos
incluyen~Bitcoin~y~Ethereum.

Por otro lado, las~\textbf{\emph{blockchains}}
\emph{\textbf{permissioned}}~restringen el acceso a un conjunto
predefinido de participantes autorizados. En este contexto, los nodos
validadores son conocidos y controlados, lo que permite optimizar el
rendimiento, reducir el coste computacional y facilitar el cumplimiento
de requisitos regulatorios. Sin embargo, este modelo introduce un mayor
grado de centralización. Un ejemplo ampliamente adoptado en entornos
empresariales es~Hyperledger Fabric.

En términos generales, las redes públicas priorizan la transparencia y
la descentralización, mientras que las redes privadas favorecen la
eficiencia operativa, el control y la privacidad, lo que explica su
adopción en contextos corporativos.~{[}2{]}

\subsubsection{2.1.2. ~Consenso entre nodos}\label{consenso-entre-nodos}

Uno de los problemas fundamentales en sistemas distribuidos consiste en
lograr que un conjunto de nodos alcance un acuerdo sobre el estado del
sistema, incluso en presencia de fallos o comportamientos maliciosos.
Este problema ha sido ampliamente estudiado en la literatura, destacando
el trabajo de~Leslie Lamport~junto con Shostak y Pease, quienes
formalizaron el consenso en entornos con fallos arbitrarios mediante
el~\emph{Byzantine Generals Problem}, sentando las bases de la
tolerancia a fallos bizantinos (\emph{Byzantine Fault Tolerance}, BFT).
{[}3{]}

En este contexto, un sistema tolerante a fallos bizantinos debe ser
capaz de alcanzar consenso incluso cuando algunos nodos presentan
comportamientos maliciosos o arbitrarios, como el envío de información
falsa o la manipulación deliberada del protocolo. Este modelo resulta
especialmente relevante en redes abiertas y sin permisos, como las
blockchain públicas, donde no existe una relación de confianza previa
entre los participantes.

No obstante, no todos los algoritmos de consenso contemplan este modelo
adversarial. Protocolos clásicos como~\emph{Paxos Protocol}~están
diseñados para entornos con fallos por parada (\emph{crash faults}), en
los que los nodos pueden fallar o dejar de responder, pero no actúan de
forma maliciosa. Esta limitación los hace inadecuados para escenarios
abiertos como blockchain, donde se asume la posible existencia de
actores maliciosos. {[}4{]}

A partir de estos fundamentos, las redes blockchain han desarrollado
mecanismos de consenso específicos que permiten mantener la coherencia
del sistema en entornos descentralizados y potencialmente hostiles.
{[}2{]}Entre los más relevantes destacan:

\begin{itemize}
\item
  \textbf{Proof of Work (PoW)}: basado en la resolución de problemas
  criptográficos que requieren un elevado coste computacional. Este
  mecanismo, utilizado por Bitcoin, proporciona seguridad al hacer
  económicamente inviable la manipulación de la cadena, aunque presenta
  limitaciones en términos de eficiencia energética y escalabilidad.
\item
  \textbf{Proof of Stake (PoS)}: sustituye el esfuerzo computacional por
  la participación económica, donde los validadores son seleccionados en
  función de la cantidad de activos bloqueados (\emph{stake}). Este
  enfoque mejora la eficiencia energética, pero introduce nuevos
  riesgos, como la concentración de poder o ataques derivados del
  problema~\emph{nothing-at-stake}.
\item
  \textbf{Delegated Proof of Stake (DPoS)}: variante de PoS en la que
  los validadores son elegidos mediante mecanismos de votación. Aumenta
  el rendimiento del sistema, aunque reduce el grado de
  descentralización.
\item
  \textbf{Proof of Authority (PoA)}: modelo en el que un conjunto
  limitado de nodos autorizados valida las transacciones. Se utiliza
  principalmente en blockchains privadas, donde la identidad de los
  participantes es conocida.
\end{itemize}

Cada uno de estos mecanismos implica distintos compromisos entre
seguridad, descentralización y rendimiento, lo que en la literatura se
describe frecuentemente como un conjunto de~\emph{trade-offs}~inherentes
al diseño de sistemas blockchain, comúnmente denominado el ``trilema de
blockchain''. {[}5{]}, {[}6{]}

Desde una perspectiva de ciberseguridad, la elección del protocolo de
consenso no solo condiciona la arquitectura del sistema, sino también su
superficie de ataque. En particular, estos compromisos influyen
directamente en la viabilidad de vectores de explotación como ataques
del 51\%, censura de transacciones o manipulación del orden de ejecución
(\emph{front-running}), especialmente en entornos abiertos y sin
permisos.

\subsubsection{2.1.3. ~Fundamentos de contratos
inteligentes}\label{fundamentos-de-contratos-inteligentes}

Los \textbf{contratos inteligentes} (\emph{smart contracts}) constituyen
una extensión funcional de la tecnología blockchain que permite la
ejecución de lógica programable de forma descentralizada. El concepto
fue introducido por~Nick Szabo~en 1994, quien los definió como ``un
protocolo de transacción informatizado que ejecuta los términos de un
contrato'' {[}7{]}

En la actualidad, los contratos inteligentes se materializan como
programas desplegados en una red blockchain que encapsulan tanto lógica
como estado persistente. Su ejecución es llevada a cabo por los nodos de
la red, los cuales deben alcanzar un resultado determinista e idéntico
para garantizar la coherencia global del sistema. Este requisito impone
restricciones relevantes en el diseño del software, especialmente en lo
relativo al uso de fuentes externas de información o funciones no
deterministas.

Cabe destacar que no todas las plataformas blockchain soportan contratos
inteligentes. Redes como~Bitcoin~incorporan un lenguaje de scripting
limitado, diseñado para priorizar la seguridad y la simplicidad,
mientras que otras plataformas como~Ethereum~permiten la ejecución de
código arbitrario, habilitando así el desarrollo de aplicaciones
complejas sobre la blockchain. {[}8{]}

\subsubsection{2.1.4. Ethereum Virtual
Machine}\label{ethereum-virtual-machine}

La adopción generalizada de contratos inteligentes se consolidó con la
aparición de~Ethereum, que introduce un entorno de ejecución específico
denominado~\textbf{\emph{Ethereum Virtual Machine}~(EVM)}. Esta máquina
virtual permite la ejecución de código de propósito general dentro de un
entorno distribuido y determinista.

El funcionamiento interno de Ethereum se describe formalmente en
el~\emph{Yellow Paper}, elaborado por~Gavin Wood, donde se definen:

\begin{itemize}
\tightlist
\item
  La arquitectura y semántica de la EVM
\item
  El modelo de ejecución de contratos y transacciones
\item
  El sistema de gas como mecanismo de control de recursos
\end{itemize}

La EVM es una máquina virtual basada en pila (\emph{stack-based}), donde
cada operación tiene un coste asociado en gas. Este mecanismo permite
limitar el consumo de recursos computacionales y actúa como protección
frente a ataques de denegación de servicio, al imponer un coste
económico a la ejecución de operaciones complejas o potencialmente
abusivas. {[}9{]}

\subsubsection{2.1.5. Propiedades de los contratos
inteligentes}\label{propiedades-de-los-contratos-inteligentes}

Desde un punto de vista técnico, los contratos inteligentes presentan
una serie de propiedades fundamentales{[}2{]}:

\begin{itemize}
\tightlist
\item
  \textbf{Determinismo}: la ejecución debe producir el mismo resultado
  en todos los nodos, garantizando la coherencia del sistema
  distribuido.
\item
  \textbf{Inmutabilidad}: una vez desplegados, los contratos no pueden
  modificarse directamente, lo que dificulta la corrección de errores.
\item
  \textbf{Transparencia}: en redes públicas, el código y el estado del
  contrato son accesibles, favoreciendo la auditabilidad.
\item
  \textbf{Ejecución descentralizada}: no existe un punto único de
  control, eliminando dependencias de entidades centrales.
\end{itemize}

\subsubsection{2.1.6. Limitaciones y
riesgos}\label{limitaciones-y-riesgos}

A pesar de sus ventajas, los contratos inteligentes presentan
importantes desafíos desde el punto de vista de la ciberseguridad
{[}2{]}:

\begin{itemize}
\tightlist
\item
  \textbf{Vulnerabilidades en el código}: errores de implementación
  pueden derivar en fallos críticos explotables.
\item
  \textbf{Problemas de control de acceso}: una gestión incorrecta de
  permisos puede permitir operaciones no autorizadas.
\item
  \textbf{Dependencia de oráculos externos}: introduce confianza en
  terceros y posibles vectores de manipulación.
\item
  \textbf{Inmutabilidad del código}: dificulta la corrección de
  vulnerabilidades tras el despliegue.
\item
  \textbf{Costes de ejecución}: el modelo de gas puede ser explotado
  para provocar condiciones de denegación de servicio.
\end{itemize}

Estas características hacen que los errores en contratos inteligentes
puedan tener consecuencias críticas, como la pérdida irreversible de
activos digitales. {[}2{]}, {[}10{]}

\subsubsection{2.1.7. Aplicaciones descentralizadas
(DApps)}\label{aplicaciones-descentralizadas-dapps}

Sobre la base del modelo de ejecución, las propiedades y las
limitaciones descritas, los contratos inteligentes permiten la
construcción de aplicaciones descentralizadas (\emph{Decentralized
Applications}, DApps), que constituyen sistemas completos en los que la
lógica de negocio se implementa parcial o totalmente mediante contratos
inteligentes. {[}11{]}

A diferencia de las aplicaciones tradicionales basadas en arquitecturas
cliente-servidor, las DApps distribuyen la lógica entre múltiples nodos
de la red, eliminando intermediarios y reduciendo la dependencia de
entidades centralizadas. Este cambio implica una redefinición del modelo
de confianza, donde los usuarios depositan su confianza en el código
ejecutado en la blockchain y en las garantías del protocolo subyacente.

Desde la perspectiva de la ciberseguridad, este paradigma introduce
implicaciones relevantes. La transparencia del código facilita su
auditoría, pero también permite a actores maliciosos analizarlo en busca
de vulnerabilidades. Asimismo, la composición de múltiples contratos y
la dependencia de servicios externos amplían la superficie de ataque,
pudiendo dar lugar a vulnerabilidades complejas en sistemas
interconectados, como ocurre en el ecosistema DeFi.

Además, la inmutabilidad de los contratos implica que los errores no
pueden corregirse fácilmente tras su despliegue, lo que incrementa el
impacto potencial de cualquier vulnerabilidad explotada. {[}12{]}

\input{capitulos/2. Estado del arte/2.03.Interoperabilidad blockchain}
\subsection{2.4. Vulnerabilidades de seguridad en contratos inteligentes
de
Ethereum}\label{vulnerabilidades-de-seguridad-en-contratos-inteligentes-de-ethereum}

Los contratos inteligentes en Ethereum presentan un modelo de seguridad
distinto al del software tradicional. Su ejecución es pública, el código
es accesible y cualquier error puede ser analizado y explotado por
actores con incentivos económicos directos. Además, la inmutabilidad
dificulta la corrección de fallos una vez desplegado el contrato, y la
interacción con otros contratos y servicios externos incrementa la
superficie de ataque. {[}10{]}, {[}30{]}

Por este motivo, una vulnerabilidad en Solidity no debe entenderse solo
como un error de programación, sino como una debilidad explotable en un
entorno adversarial. En la práctica, los problemas de seguridad no se
limitan a fallos técnicos, sino que también incluyen errores de diseño y
de lógica de negocio. {[}31{]}

Para su análisis, resulta útil agrupar las vulnerabilidades en cuatro
categorías principales:

\begin{itemize}
\tightlist
\item
  Vulnerabilidades técnicas de ejecución
\item
  Vulnerabilidades de control y privilegios
\item
  Vulnerabilidades económicas y dependencia del entorno
\item
  Errores lógicos de negocio
\end{itemize}

\subsubsection{2.4.1. Vulnerabilidades técnicas de
ejecución}\label{vulnerabilidades-tuxe9cnicas-de-ejecuciuxf3n}

Este grupo incluye errores directamente relacionados con la ejecución en
la EVM y la implementación en Solidity.

Una de las más conocidas es la~\textbf{reentrancy}, que ocurre cuando un
contrato realiza una llamada externa antes de actualizar su estado
interno. Esto permite que el contrato receptor vuelva a ejecutar la
función vulnerable con un estado inconsistente. A pesar de que existen
patrones de mitigación bien conocidos, sigue siendo una de las
vulnerabilidades más relevantes en Ethereum. {[}32{]}

Otra categoría importante son los~\textbf{problemas aritméticos}, como
overflow y underflow. Aunque Solidity 0.8 introduce comprobaciones
automáticas, estos errores siguen siendo relevantes en contratos
antiguos, en bloques~unchecked~o en cálculos financieros incorrectos.

También destacan las~\textbf{llamadas externas inseguras}, ya que
cualquier~call~transfiere el control de ejecución a otro contrato. Esto
puede provocar comportamientos inesperados o ejecución de lógica
maliciosa. En esta línea, el uso de~delegatecall~es especialmente
crítico, ya que ejecuta código externo sobre el almacenamiento del
contrato llamador, pudiendo comprometer completamente su estado.~
{[}30{]}

Por último, los~\textbf{ataques de denegación de servicio (DoS)}~son
frecuentes en este entorno. No buscan tumbar el sistema, sino bloquear
funciones críticas, por ejemplo, mediante bucles no acotados o
reversiones provocadas por contratos externos. {[}30{]}

\subsubsection{2.4.2. Vulnerabilidades de control y
privilegios}\label{vulnerabilidades-de-control-y-privilegios}

Muchas vulnerabilidades reales no se deben a errores técnicos complejos,
sino a problemas en el control de acceso.

Esto incluye funciones sensibles sin restricciones, roles mal definidos
o privilegios excesivos. Un error típico es permitir que cualquier
usuario ejecute funciones críticas como transferencias de fondos o
cambios de configuración. {[}33{]}

Un caso especialmente problemático es el uso de~tx.origin~para
autenticación. Este valor representa el origen externo de la
transacción, no el llamador inmediato, lo que permite ataques mediante
contratos intermedios.

También son relevantes los problemas en contratos upgradeables,
especialmente los relacionados con inicializadores. Si una
función~initialize~no está correctamente protegida, un atacante puede
ejecutarla y asumir el control del contrato. {[}34{]}

\subsubsection{2.4.3. Vulnerabilidades económicas y dependencia del
entorno}\label{vulnerabilidades-econuxf3micas-y-dependencia-del-entorno}

En muchos casos, el problema no está en el código en sí, sino en cómo
interactúa el contrato con su entorno.

El ejemplo más claro es el~\textbf{front-running}, donde un atacante
observa transacciones en la mempool y envía otras con mayor prioridad
para ejecutarse antes. Este fenómeno forma parte del problema más amplio
del~\textbf{MEV (Maximal Extractable Value)}~y afecta especialmente a
aplicaciones DeFi. {[}35{]}

Otro riesgo importante es la~\textbf{manipulación de oráculos}. Si un
contrato depende de precios externos que pueden alterarse temporalmente,
un atacante puede aprovechar esa situación para obtener beneficios
indebidos, por ejemplo, en protocolos de préstamo. {[}33{]}

En esta categoría también se incluyen los problemas
de~\textbf{aleatoriedad insegura}~y el uso de~block.timestamp~como
fuente de decisiones críticas. Estas variables no son completamente
fiables ni impredecibles, por lo que pueden ser manipuladas o
anticipadas en determinados escenarios.

\subsubsection{2.4.4. Errores lógicos de
negocio}\label{errores-luxf3gicos-de-negocio}

Los errores más difíciles de detectar son los relacionados con la lógica
del contrato.

En estos casos, el código puede ser correcto desde el punto de vista
técnico, pero implementar una lógica incorrecta. Esto incluye errores en
cálculos de balances, distribución de recompensas, gestión de estados o
validación de condiciones. {[}31{]}

Este tipo de vulnerabilidades es especialmente relevante porque suele
escapar a las herramientas automáticas. Además, muchos ataques reales
combinan errores lógicos con otros factores, como manipulación de
precios o condiciones de ejecución. {[}35{]}

\subsubsection{2.4.5. Conclusión}\label{conclusiuxf3n}

En conjunto, la seguridad en contratos inteligentes no puede abordarse
únicamente mediante la detección de patrones conocidos. Aunque
vulnerabilidades como reentrancy o overflow siguen siendo relevantes,
los problemas más complejos suelen estar relacionados con el diseño del
sistema, la interacción con otros componentes y la lógica de negocio.

Esto implica que una auditoría efectiva debe analizar no solo el código,
sino también su comportamiento, sus dependencias y los incentivos que lo
rodean.

\subsection{2.5. Ataques Reales}\label{ataques-reales}

El estudio de incidentes reales constituye una fuente de conocimiento
imprescindible en el ámbito de la seguridad de contratos inteligentes. A
diferencia de los entornos de prueba, los ataques producidos en redes de
producción demuestran cómo las vulnerabilidades teóricas se traducen en
pérdidas económicas concretas e irreversibles.

En esta sección se analizan cuatro casos históricos representativos,
seleccionados por su relevancia técnica y su correspondencia directa con
cada una de las categorías de vulnerabilidades descritas en la sección
anterior (4.4. Vulnerabilidades de seguridad en contratos inteligentes
de Ethereum).

\subsubsection{2.5.1. Vulnerabilidad técnica de ejecución: The DAO
(2016)}\label{vulnerabilidad-tuxe9cnica-de-ejecuciuxf3n-the-dao-2016}

El ataque contra The DAO, producido en junio de 2016, representa el caso
más paradigmático de explotación de una vulnerabilidad de
\emph{reentrancy} en la historia de Ethereum {[}36{]}. The DAO era un
fondo de inversión descentralizado desplegado sobre Ethereum que había
recaudado aproximadamente 150 millones de dólares en ETH procedentes de
más de once mil participantes, convirtiéndose en uno de los mayores
proyectos de financiación colectiva hasta ese momento {[}37{]}.

La vulnerabilidad residía en la función withdraw() del contrato. El
flujo de ejecución enviaba los fondos en ETH al destinatario antes de
actualizar el balance interno correspondiente. Esta secuencia incorrecta
permitió a un atacante desplegar un contrato receptor cuya función de
\emph{fallback} re-invocaba withdraw() de forma recursiva, extrayendo
fondos de manera repetida mientras el saldo del contrato permanecía sin
modificar. El ataque se ejecutó en múltiples iteraciones dentro de una
misma cadena de llamadas, drenando aproximadamente 3,6 millones de ETH
{[}38{]}.

Las consecuencias del incidente trascendieron lo puramente económico. La
comunidad de Ethereum debatió durante semanas sobre cómo dar respuesta
al ataque, dado el carácter inmutable del protocolo. Finalmente, se
ejecutó un \emph{hard fork} en el bloque 1.920.000, el 20 de julio de
2016, que revirtió el estado de la cadena para recuperar los fondos. Una
fracción de la comunidad rechazó esta intervención por considerarla
contraria a los principios de inmutabilidad de la tecnología blockchain,
lo que dio lugar a la bifurcación conocida como Ethereum Classic
{[}37{]}.

Este caso puso de manifiesto que el patrón de diseño correcto para
gestionar transferencias es el denominado
\emph{checks-effects-interactions}: primero verificar condiciones,
después actualizar el estado del contrato y, por último, realizar
cualquier llamada externa. Su incumplimiento en The DAO, a pesar de la
aparente sencillez del principio, resultó en pérdidas valoradas en torno
a 60 millones de dólares al precio del momento {[}36{]}.

\subsubsection{2.5.2. Vulnerabilidad de control y privilegios: Poly
Network
(2021)}\label{vulnerabilidad-de-control-y-privilegios-poly-network-2021}

El ataque a Poly Network, ocurrido el 10 de agosto de 2021, ilustra con
especial claridad las consecuencias de un control de acceso deficiente
en contratos que gestionan operaciones críticas de alto valor {[}39{]}.
Poly Network es un protocolo de interoperabilidad \emph{cross-chain} que
conecta diversas redes blockchain, lo que amplía considerablemente su
superficie de ataque.

El vector de explotación se basó en una ausencia de restricciones en la
cadena de llamadas entre contratos del protocolo. La función
verifyHeaderAndExecuteTx(), presente en el contrato
EthCrossChainManager, permitía invocar internamente la función
PutCurEpochConPubKeyBytes() del contrato EthCrossChainData, encargada de
actualizar la clave pública del grupo de \emph{keepers} con permisos
para ejecutar transacciones \emph{cross-chain}. Al no existir ninguna
comprobación que restringiese qué direcciones podían realizar esta
invocación de forma indirecta, el atacante fue capaz de sustituir los
\emph{keepers} legítimos por una dirección bajo su propio control.

La operación se replicó de forma simultánea en tres redes: Ethereum, BNB
Chain y Polygon; lo que supuso un control inmediato sobre la práctica
totalidad de los fondos bloqueados en el protocolo, con un valor
estimado de 611 millones de dólares, convirtiéndose en el mayor robo de
la historia del ecosistema DeFi hasta esa fecha.

El desenlace del incidente fue inusual: el atacante devolvió la
totalidad de los fondos en los días posteriores, argumentando haber
actuado con el propósito de demostrar la vulnerabilidad. No obstante, el
caso evidencia que errores de diseño en el modelo de permisos
(independientemente de su recuperabilidad) pueden comprometer
completamente la integridad de un protocolo. Una correcta implementación
de control de acceso, mediante modificadores de autorización explícitos
y separación de privilegios entre contratos, habría impedido la escalada
de permisos aprovechada en este ataque {[}40{]}.

\subsubsection{2.5.3. Vulnerabilidad económica y dependencia del
entorno: bZx
(2020)}\label{vulnerabilidad-econuxf3mica-y-dependencia-del-entorno-bzx-2020}

Los ataques contra el protocolo bZx, producidos en febrero de 2020,
constituyen uno de los primeros ejemplos documentados de explotación
combinada de \emph{flash loans} y manipulación de oráculos de precio a
escala en el ecosistema DeFi. Su relevancia radica no solo en las
pérdidas económicas, sino en haber demostrado que la composabilidad de
los protocolos DeFi puede convertirse en un vector de ataque cuando los
componentes individuales no contemplan escenarios de manipulación
externa {[}41{]}.

En el primero de los dos incidentes, acaecido el 14 de febrero de 2020,
el atacante obtuvo un préstamo flash de 10.000 ETH a través del
protocolo dYdX sin necesidad de aportar colateral alguno. Con una parte
de estos fondos abrió simultáneamente una posición corta apalancada
sobre el par WBTC/ETH en bZx y utilizó el volumen restante para adquirir
WBTC en Uniswap, manipulando artificialmente el precio de mercado del
par en ese pool. El protocolo bZx, al depender de Uniswap como fuente de
precios en tiempo real, calculó incorrectamente la valoración de la
posición del atacante, quien la cerró obteniendo un beneficio aproximado
de 350.000 dólares una vez devuelto el préstamo flash, todo dentro de
una única transacción {[}42{]}.

El segundo ataque, cuatro días después, siguió una mecánica similar pero
orientada a la manipulación del precio del token sUSD en Kyber Network.
Al inflar artificialmente su cotización, el atacante logró depositar
colateral sobrevalorado en bZx y extraer en préstamo un volumen de
activos muy superior al que le habría correspondido según condiciones de
mercado normales {[}42{]}.

Ambos ataques pusieron de manifiesto que la seguridad de un protocolo
DeFi no puede analizarse de forma aislada: la dependencia de precios
obtenidos de fuentes únicas y manipulables en tiempo real constituye una
vulnerabilidad estructural. La adopción de oráculos descentralizados con
agregación de múltiples fuentes de precio (como los proporcionados por
protocolos del tipo Chainlink) y el uso de precios promediados en el
tiempo (TWAP, \emph{Time-Weighted Average Price}) son las principales
contramedidas frente a este patrón de ataque.

\subsubsection{2.5.4. Error lógico de negocio: Euler Finance
(2023)}\label{error-luxf3gico-de-negocio-euler-finance-2023}

El ataque a Euler Finance, ejecutado el 13 de marzo de 2023, representa
un ejemplo de elevada complejidad técnica dentro de la categoría de
errores lógicos, precisamente porque la vulnerabilidad no residía en un
patrón de codificación inseguro clásico, sino en una interacción no
prevista entre mecanismos financieros del propio protocolo {[}43{]}.

Euler Finance era un protocolo de préstamos descentralizado en Ethereum
que introducía el concepto de \emph{sub-cuentas} para gestionar
posiciones de colateral y deuda de forma independiente. La
vulnerabilidad se encontraba en la función donateToReserves(),
incorporada en una actualización anterior. Esta función permitía a un
usuario transferir sus propios activos a las reservas del protocolo,
pero no verificaba que dicha operación no dejase la cuenta del donante
en un estado de insolvencia. En condiciones normales, esta situación
habría sido detectada por el sistema de liquidación, pero el atacante
aprovechó el mecanismo de autoliquidación del protocolo de una forma no
contemplada en su diseño {[}44{]}.

El flujo del ataque comenzó con la obtención de un préstamo flash de 30
millones de DAI desde el protocolo Aave. Con estos fondos, el atacante
depositó colateral en Euler y, mediante operaciones de apalancamiento
repetidas, construyó posiciones de aproximadamente 195 millones de eDAI
(depósitos) y 200 millones de dDAI (deuda). A continuación, utilizó
donateToReserves() para crear intencionalmente una posición insolvente y
acto seguido activó la autoliquidación. El sistema de Euler, al calcular
el bonus de liquidación sobre una posición de ese tamaño, transfirió al
atacante los fondos de las reservas del protocolo en una cuantía muy
superior a la deuda original {[}45{]}.

La pérdida total ascendió a aproximadamente 197 millones de dólares. Sin
embargo, el atacante devolvió la práctica totalidad de los fondos
semanas después, tras una negociación directa con el equipo de Euler
mediante mensajes publicados en la cadena de bloques. El incidente
ilustra cómo la ausencia de validaciones sobre invariantes de negocio
(en este caso, que ninguna operación debería permitir dejar una cuenta
en estado de insolvencia) puede generar vulnerabilidades que escapan
tanto a la revisión manual del código como a las herramientas de
análisis estático-convencionales {[}44{]}.

\input{capitulos/2. Estado del arte/2.06.Herramientas de Análisis}
\input{capitulos/2. Estado del arte/2.07.Sintesis}

\chapter{Objetivos concretos y metodología de trabajo}
\input{capitulos/3. Objetivos concretos y metodología de trabajo}

\chapter{Desarrollo específico de la contribución}
\input{capitulos/4. Desarrollo específico de la contribución/4.00.Introducción}
\input{capitulos/4. Desarrollo específico de la contribución/4.01.Requisitos}
\input{capitulos/4. Desarrollo específico de la contribución/4.02.Arquitectura}
\section{4.3.Runner}\label{runner}

\subsection{4.3. Módulo Runner}\label{muxf3dulo-runner}

El módulo Runner constituye la capa encargada de la interacción con las
herramientas externas utilizadas durante el proceso de análisis. Su
principal responsabilidad consiste en gestionar la ejecución de Slither,
Mythril y Echidna, así como preparar el entorno necesario para
garantizar la correcta realización del análisis.

La existencia de un módulo específico para esta tarea permite desacoplar
la lógica propia de EVMAudit de las particularidades de cada
herramienta, facilitando la incorporación de nuevas soluciones de
análisis en futuras versiones.

!TODO: insertar referencia a la Figura X (diagrama de arquitectura).

\subsubsection{4.3.1. Responsabilidades del
módulo}\label{responsabilidades-del-muxf3dulo}

Las principales responsabilidades implementadas por el módulo Runner son
las siguientes:

\begin{itemize}
\tightlist
\item
  detección automática del nombre del contrato
\item
  configuración de la versión del compilador Solidity
\item
  ejecución de Slither
\item
  ejecución de Mythril
\item
  ejecución de Echidna
\item
  almacenamiento de resultados intermedios
\item
  gestión de errores producidos durante la ejecución.
\end{itemize}

Gracias a esta aproximación, el resto de módulos del sistema pueden
trabajar de forma independiente sin necesidad de conocer los detalles
específicos de cada herramienta externa.

\subsubsection{4.3.2. Detección automática del nombre del
contrato}\label{detecciuxf3n-automuxe1tica-del-nombre-del-contrato}

Durante las primeras fases de desarrollo se observó que el nombre del
fichero Solidity no siempre coincide con el nombre del contrato definido
en su interior.

Por este motivo se implementó un mecanismo de detección automática que
analiza el código fuente y extrae el nombre real del contrato utilizando
expresiones regulares.

Esta decisión permite evitar errores durante las fases posteriores del
análisis y elimina la necesidad de que el usuario proporcione
manualmente dicha información.

Además, este mecanismo resulta especialmente útil cuando la librería es
utilizada desde otras aplicaciones externas, ya que reduce la cantidad
de parámetros que deben proporcionarse.

!TODO: insertar fragmento de código de
\texttt{detect\_contract\_name()}.

\subsubsection{4.3.3. Configuración automática del
compilador}\label{configuraciuxf3n-automuxe1tica-del-compilador}

Las herramientas de análisis empleadas dependen de la versión del
compilador Solidity utilizada por el contrato.

Con el objetivo de garantizar la reproducibilidad del análisis y evitar
incompatibilidades entre versiones, EVMAudit analiza automáticamente la
directiva \texttt{pragma\ solidity} presente en el contrato y configura
la versión correspondiente mediante la utilidad \texttt{solc-select}.

Este mecanismo permite adaptar dinámicamente el entorno de ejecución y
facilita el análisis de contratos desarrollados con diferentes versiones
del lenguaje.

!TODO: insertar referencia cruzada a la Sección 2.X (Solidity y
compilador).

!TODO: insertar fragmento de código de \texttt{\_set\_solc\_version()}.

\subsubsection{4.3.4. Ejecución de
Slither}\label{ejecuciuxf3n-de-slither}

Slither constituye la primera herramienta utilizada durante el proceso
de análisis debido a su rapidez y a la gran cantidad de detectores
disponibles.

La función encargada de su ejecución verifica previamente que la
herramienta se encuentre instalada y posteriormente lanza el análisis
sobre el contrato proporcionado.

Una vez finalizada la ejecución, la salida obtenida se almacena en
formato JSON para garantizar la trazabilidad de los resultados.

Durante la implementación fue necesario contemplar ciertos
comportamientos específicos de Slither. En particular, la herramienta
devuelve el código de salida 255 cuando detecta vulnerabilidades, aunque
dicho comportamiento no representa un error real.

Por este motivo, EVMAudit incorpora una lógica específica para
interpretar correctamente este caso y continuar con el proceso de
análisis.

!TODO: insertar fragmento de código de \texttt{run\_slither()}.

!TODO: referencia cruzada a la Sección 2.X (Slither).

\subsubsection{4.3.5. Ejecución de
Mythril}\label{ejecuciuxf3n-de-mythril}

La segunda fase del análisis se realiza mediante Mythril, herramienta
basada en ejecución simbólica.

A diferencia de Slither, Mythril presenta tiempos de ejecución
superiores y requiere parámetros adicionales relacionados con la
profundidad máxima de exploración y el tiempo máximo permitido para
completar el análisis.

Durante el desarrollo se incorporaron mecanismos para:

\begin{itemize}
\tightlist
\item
  limitar el tiempo máximo de ejecución
\item
  controlar posibles errores durante el análisis
\item
  procesar las salidas generadas por la herramienta
\item
  preservar los resultados originales.
\end{itemize}

Estas medidas permiten evitar que un contrato especialmente complejo
bloquee el proceso completo de análisis.

!TODO: insertar fragmento de código de \texttt{run\_mythril()}.

!TODO: referencia cruzada a la Sección 2.X (Mythril).

\subsubsection{4.3.6. Ejecución de
Echidna}\label{ejecuciuxf3n-de-echidna}

Una vez generadas las propiedades correspondientes, EVMAudit ejecuta
Echidna para realizar una fase adicional de validación mediante fuzzing.

La integración de Echidna presentó una complejidad superior a la
observada en Slither y Mythril debido a ciertas particularidades de la
herramienta y a las diferencias existentes entre versiones.

Durante el desarrollo fue necesario incorporar mecanismos específicos
para:

\begin{itemize}
\tightlist
\item
  reconstruir nombres de propiedades
\item
  interpretar adecuadamente las salidas generadas
\item
  distinguir entre pruebas superadas y vulnerabilidades explotables
\item
  gestionar posibles inconsistencias en los resultados.
\end{itemize}

Estas adaptaciones permitieron integrar Echidna dentro del flujo general
de EVMAudit sin alterar el resto de componentes de la arquitectura.

!TODO: insertar fragmento de código de \texttt{run\_echidna()}.

!TODO: referencia cruzada a la Sección 2.X (Echidna).

\subsubsection{4.3.7. Persistencia de resultados
intermedios}\label{persistencia-de-resultados-intermedios}

Con el fin de favorecer la trazabilidad y la reproducibilidad del
análisis, EVMAudit almacena las salidas originales generadas por las
distintas herramientas utilizadas.

Esta decisión permite:

\begin{itemize}
\tightlist
\item
  inspeccionar manualmente los resultados obtenidos
\item
  reproducir análisis posteriores
\item
  depurar posibles errores
\item
  conservar evidencias originales.
\end{itemize}

La preservación de las salidas originales resulta especialmente
relevante en el contexto de auditorías de seguridad, donde la
trazabilidad de los hallazgos constituye un aspecto fundamental.

\section{4.4.Normalizer}\label{normalizer}

\subsection{4.4. Módulo de
normalización}\label{muxf3dulo-de-normalizaciuxf3n}

Uno de los principales problemas identificados durante el estudio de las
herramientas existentes fue la heterogeneidad de los formatos utilizados
para representar las vulnerabilidades detectadas.

Slither y Mythril generan estructuras JSON completamente diferentes,
tanto en la forma de representar la severidad como en la descripción de
los hallazgos o la identificación de las localizaciones afectadas.

Esta situación dificulta la comparación directa de resultados y hace
necesaria una fase previa de transformación antes de poder aplicar
mecanismos de correlación.

Para resolver este problema se desarrolló un módulo de normalización
encargado de transformar las salidas heterogéneas en una estructura
común.

\subsubsection{4.4.1. Objetivos del proceso de
normalización}\label{objetivos-del-proceso-de-normalizaciuxf3n}

El proceso de normalización persigue varios objetivos:

\begin{itemize}
\tightlist
\item
  unificar la representación de vulnerabilidades;
\item
  facilitar la correlación entre herramientas;
\item
  preservar las evidencias originales;
\item
  simplificar la generación posterior de informes;
\item
  desacoplar el resto de módulos de las particularidades de cada
  herramienta.
\end{itemize}

Gracias a ello, los componentes posteriores pueden operar sobre una
estructura homogénea independientemente del origen de los resultados.

\subsubsection{4.4.2. Estructura común de los
hallazgos}\label{estructura-comuxfan-de-los-hallazgos}

Cada vulnerabilidad normalizada se representa mediante un conjunto común
de atributos.

Entre los campos principales se encuentran:

\begin{itemize}
\tightlist
\item
  título;
\item
  descripción;
\item
  severidad;
\item
  categoría;
\item
  contrato afectado;
\item
  función afectada;
\item
  localización;
\item
  identificador SWC;
\item
  herramienta de origen;
\item
  evidencia original.
\end{itemize}

La inclusión de la salida original permite mantener la trazabilidad del
análisis y facilita la revisión manual de los resultados.

!TODO: insertar diagrama UML del objeto Finding.

!TODO: revisar nombres exactos de los atributos.

\subsubsection{4.4.3. Normalización de resultados de
Slither}\label{normalizaciuxf3n-de-resultados-de-slither}

La función de normalización correspondiente a Slither transforma la
estructura JSON generada por la herramienta y extrae la información
relevante necesaria para construir los hallazgos comunes utilizados por
EVMAudit.

Durante esta fase se realiza además la asociación de los detectores
propios de Slither con los identificadores SWC correspondientes.

Este proceso permite traducir la terminología específica de Slither a
una representación independiente de la herramienta utilizada.

!TODO: insertar fragmento de código de
\texttt{normalize\_slither\_output()}.

\subsubsection{4.4.4. Normalización de resultados de
Mythril}\label{normalizaciuxf3n-de-resultados-de-mythril}

De forma análoga, la salida generada por Mythril es procesada para
obtener una representación equivalente a la utilizada por Slither.

La transformación permite unificar:

\begin{itemize}
\tightlist
\item
  niveles de severidad;
\item
  nombres de vulnerabilidades;
\item
  identificadores SWC;
\item
  información contextual del hallazgo.
\end{itemize}

Gracias a ello, ambas herramientas pueden ser tratadas posteriormente de
forma transparente por el módulo de correlación.

!TODO: insertar fragmento de código de
\texttt{normalize\_mythril\_output()}.

\subsubsection{4.4.5. Clasificación de
vulnerabilidades}\label{clasificaciuxf3n-de-vulnerabilidades}

Además de unificar la representación de los hallazgos, el módulo de
normalización clasifica las vulnerabilidades detectadas siguiendo las
categorías definidas durante el estudio del estado del arte.

Entre las categorías utilizadas se encuentran:

\begin{itemize}
\tightlist
\item
  vulnerabilidades de control de acceso;
\item
  vulnerabilidades económicas;
\item
  vulnerabilidades relacionadas con la lógica de negocio;
\item
  vulnerabilidades asociadas a la ejecución del contrato.
\end{itemize}

Esta clasificación permite mantener la coherencia entre la parte teórica
del trabajo y la implementación desarrollada.

!TODO: referencia cruzada a la Sección 2.X (clasificación de
vulnerabilidades).

\subsubsection{4.4.6. Preservación de
evidencias}\label{preservaciuxf3n-de-evidencias}

Uno de los principios adoptados durante el diseño de EVMAudit fue evitar
la pérdida de información durante las fases intermedias del análisis.

Por este motivo, cada hallazgo normalizado conserva una referencia a la
salida original generada por la herramienta correspondiente.

Esta decisión facilita:

\begin{itemize}
\tightlist
\item
  la revisión manual por parte del auditor;
\item
  la trazabilidad del proceso;
\item
  la depuración de errores;
\item
  la reproducibilidad de los resultados.
\end{itemize}

La conservación de evidencias resulta especialmente importante en
entornos de auditoría y análisis forense, donde es necesario justificar
el origen de cada vulnerabilidad identificada.

!TODO: insertar diagrama del proceso de normalización.

!TODO: insertar referencia a la Figura X.

\section{4.5.Correlator}\label{correlator}

\subsection{4.5. Módulo de
correlación}\label{muxf3dulo-de-correlaciuxf3n}

Una vez normalizados los resultados obtenidos por las herramientas de
análisis, EVMAudit aplica un proceso de correlación cuyo objetivo
consiste en identificar vulnerabilidades equivalentes detectadas por
distintas herramientas y agruparlas en un único hallazgo.

Este módulo constituye el núcleo de la contribución desarrollada en el
presente Trabajo Fin de Máster, ya que permite reducir la redundancia de
información y proporcionar una visión más estructurada de las
vulnerabilidades detectadas.

La necesidad de incorporar esta fase surge de una de las principales
limitaciones observadas durante el estudio del estado del arte. Aunque
herramientas como Slither y Mythril son capaces de detectar un gran
número de vulnerabilidades, cada una de ellas presenta los resultados
utilizando nomenclaturas y estructuras diferentes, generando además
múltiples duplicidades que dificultan la revisión manual por parte del
auditor.

Por este motivo, EVMAudit incorpora un mecanismo propio de correlación
encargado de consolidar la información procedente de ambas herramientas.

\subsubsection{4.5.1. Objetivos del proceso de
correlación}\label{objetivos-del-proceso-de-correlaciuxf3n}

Los objetivos perseguidos por este módulo son los siguientes:

\begin{itemize}
\tightlist
\item
  reducir la redundancia de información
\item
  agrupar vulnerabilidades equivalentes
\item
  incrementar la confianza en los hallazgos detectados
\item
  facilitar la interpretación de resultados
\item
  proporcionar una representación unificada de las vulnerabilidades.
\end{itemize}

La correlación no pretende eliminar completamente los falsos positivos
generados por las herramientas utilizadas, sino proporcionar una capa
adicional de análisis que facilite el trabajo posterior del auditor.

\subsubsection{4.5.2. Estrategia de
correlación}\label{estrategia-de-correlaciuxf3n}

El proceso implementado en EVMAudit se basa en una estrategia de
correlación mediante reglas.

Dos vulnerabilidades se consideran equivalentes cuando cumplen
simultáneamente las siguientes condiciones:

\begin{itemize}
\tightlist
\item
  afectan al mismo contrato
\item
  afectan a la misma función
\item
  poseen el mismo identificador SWC.
\end{itemize}

De este modo, la clave utilizada para la correlación puede expresarse de
la siguiente forma:

\begin{Shaded}
\begin{Highlighting}[]
\NormalTok{Contrato + Función + SWC}
\end{Highlighting}
\end{Shaded}

Esta aproximación permite agrupar vulnerabilidades que, aunque procedan
de herramientas diferentes, representan realmente el mismo problema de
seguridad.

La elección del identificador SWC como criterio común permite disponer
de una referencia independiente de la nomenclatura utilizada por cada
herramienta.

!TODO: insertar referencia cruzada a la Sección 2.X (SWC Registry).

!TODO: insertar pseudocódigo del algoritmo de correlación.

\subsubsection{4.5.3. Hallazgos confirmados y hallazgos
detectados}\label{hallazgos-confirmados-y-hallazgos-detectados}

Una vez realizado el proceso de agrupación, EVMAudit distingue entre dos
tipos de vulnerabilidades.

\paragraph{Hallazgos confirmados}\label{hallazgos-confirmados}

Se consideran confirmados aquellos hallazgos que han sido detectados por
más de una herramienta.

La coincidencia entre distintas técnicas de análisis aumenta la
confianza asociada a la vulnerabilidad y proporciona una mayor evidencia
de su existencia.

\paragraph{Hallazgos detectados}\label{hallazgos-detectados}

Corresponden a vulnerabilidades identificadas por una única herramienta.

Aunque presentan un menor nivel de confianza, siguen siendo incluidas en
el informe final con el objetivo de evitar la pérdida de información
potencialmente relevante.

Esta clasificación permite priorizar posteriormente la revisión manual
realizada por el auditor.

\subsubsection{4.5.4. Nivel de confianza}\label{nivel-de-confianza}

Con el objetivo de proporcionar información adicional sobre la
fiabilidad de cada hallazgo, EVMAudit asigna un nivel de confianza
basado en las herramientas que han detectado la vulnerabilidad.

La confianza asociada no representa una medida probabilística, sino un
indicador orientativo derivado del número de herramientas que coinciden
en el mismo hallazgo.

Este mecanismo permite distinguir aquellas vulnerabilidades respaldadas
por varias herramientas de aquellas detectadas únicamente por una de
ellas.

No obstante, la interpretación final continúa dependiendo del criterio
del auditor, ya que la correlación implementada no sustituye la
validación manual.

!TODO: revisar nomenclatura definitiva utilizada por el código
(\texttt{confidence\_score}).

\subsubsection{4.5.5. Gestión de
severidades}\label{gestiuxf3n-de-severidades}

Cuando una vulnerabilidad es detectada por varias herramientas, puede
ocurrir que cada una de ellas asigne niveles de severidad diferentes.

Para resolver esta situación, EVMAudit conserva la severidad más alta
reportada por las herramientas implicadas.

Este enfoque se adopta siguiendo un criterio conservador, priorizando la
revisión de aquellas vulnerabilidades que potencialmente puedan tener un
mayor impacto.

\subsubsection{4.5.6. Preservación de
evidencias}\label{preservaciuxf3n-de-evidencias}

Durante el proceso de correlación se mantiene la información procedente
de todas las herramientas que han participado en la detección.

De esta forma, cada hallazgo correlacionado conserva:

\begin{itemize}
\tightlist
\item
  las herramientas que lo han identificado
\item
  las evidencias originales asociadas
\item
  las localizaciones afectadas
\item
  la información contextual proporcionada por cada herramienta.
\end{itemize}

Esta decisión garantiza la trazabilidad del análisis y facilita la
revisión manual posterior.

\subsubsection{4.5.7. Beneficios del proceso de
correlación}\label{beneficios-del-proceso-de-correlaciuxf3n}

La incorporación de esta fase proporciona diversas ventajas:

\begin{itemize}
\tightlist
\item
  reducción de duplicidades
\item
  mejora de la legibilidad del informe final
\item
  aumento de la confianza en determinados hallazgos
\item
  simplificación del proceso de auditoría
\item
  independencia respecto a las herramientas utilizadas.
\end{itemize}

La correlación constituye, por tanto, uno de los principales elementos
diferenciadores de EVMAudit frente al uso individual de las herramientas
integradas.

!TODO: insertar diagrama del proceso de correlación.

!TODO: insertar ejemplo real antes y después de la correlación.

!TODO: insertar referencia a la Figura X.

\section{4.6.SWC Catalog}\label{swc-catalog}

\subsection{4.6. Catálogo de vulnerabilidades
SWC}\label{catuxe1logo-de-vulnerabilidades-swc}

Con el objetivo de desacoplar la fase de detección de vulnerabilidades
de las fases posteriores de validación y generación de pruebas, se
implementó un catálogo de vulnerabilidades basado en la clasificación
Smart Contract Weakness Classification (SWC).

Este catálogo constituye una capa intermedia que permite asociar los
hallazgos detectados con información adicional independiente de las
herramientas utilizadas.

Gracias a ello, EVMAudit puede reutilizar la información obtenida
durante el proceso de correlación y emplearla posteriormente para la
generación automática de propiedades orientadas a Echidna.

\subsubsection{4.6.1. Objetivos del
catálogo}\label{objetivos-del-catuxe1logo}

El catálogo desarrollado persigue los siguientes objetivos:

\begin{itemize}
\tightlist
\item
  unificar la representación de vulnerabilidades
\item
  desacoplar la lógica de generación de pruebas
\item
  facilitar la incorporación de nuevas plantillas
\item
  centralizar información sobre las debilidades conocidas
\item
  favorecer la mantenibilidad del sistema.
\end{itemize}

Esta aproximación permite que las fases posteriores del análisis no
dependan directamente de las particularidades de Slither o Mythril.

\subsubsection{4.6.2. Estructura del
catálogo}\label{estructura-del-catuxe1logo}

Cada entrada del catálogo contiene información asociada a una
vulnerabilidad concreta.

Entre los datos almacenados se encuentran:

\begin{itemize}
\tightlist
\item
  identificador SWC
\item
  nombre de la vulnerabilidad
\item
  descripción
\item
  severidad
\item
  plantilla asociada
\item
  limitaciones conocidas
\item
  posibilidad de validación automática.
\end{itemize}

La información almacenada permite enriquecer los hallazgos
correlacionados y proporcionar información adicional durante las fases
posteriores del análisis.

!TODO: insertar tabla con varios ejemplos del catálogo.

\subsubsection{4.6.3. Asociación entre detectores y
SWC}\label{asociaciuxf3n-entre-detectores-y-swc}

Una de las principales dificultades observadas durante el desarrollo fue
la ausencia de una correspondencia directa entre los detectores
utilizados por cada herramienta.

Para resolver este problema se definieron tablas de asociación que
permiten traducir los detectores específicos a identificadores SWC.

Gracias a esta estrategia es posible utilizar un lenguaje común durante
todo el proceso de análisis y evitar dependencias directas respecto a la
nomenclatura propia de cada herramienta.

Esta decisión resulta especialmente relevante para el módulo de
correlación descrito anteriormente.

!TODO: insertar referencia cruzada a la Sección 4.4.

\subsubsection{4.6.4. Reutilización del
catálogo}\label{reutilizaciuxf3n-del-catuxe1logo}

Además de servir como mecanismo de clasificación, el catálogo es
utilizado posteriormente durante la generación de propiedades para
Echidna.

A partir del identificador SWC asociado a cada vulnerabilidad, EVMAudit
recupera automáticamente la plantilla correspondiente y genera las
estructuras necesarias para la fase de fuzzing.

Esta aproximación permite separar claramente:

\begin{itemize}
\tightlist
\item
  detección
\item
  correlación
\item
  generación de pruebas.
\end{itemize}

La independencia entre estas fases facilita futuras ampliaciones y
simplifica el mantenimiento del sistema.

\subsubsection{4.6.5. Extensibilidad del
catálogo}\label{extensibilidad-del-catuxe1logo}

La utilización de un catálogo independiente permite incorporar nuevas
vulnerabilidades sin necesidad de modificar el resto de módulos
implementados.

De esta forma, la incorporación de nuevas entradas únicamente requiere:

\begin{enumerate}
\def\labelenumi{\arabic{enumi}.}
\tightlist
\item
  definir el identificador SWC correspondiente
\item
  añadir la información descriptiva asociada
\item
  incorporar la plantilla de generación deseada.
\end{enumerate}

Este diseño favorece la evolución futura de la herramienta y permite
adaptar EVMAudit a nuevas versiones del ecosistema Ethereum.

!TODO: insertar diagrama de relación entre SWC Catalog y Echidna
Adapter.

!TODO: insertar referencia a la Figura X.

\section{4.7.Echidna adaptar}\label{echidna-adaptar}

\subsection{4.7. Adaptador para Echidna y generación automática de
propiedades}\label{adaptador-para-echidna-y-generaciuxf3n-automuxe1tica-de-propiedades}

Una vez finalizado el proceso de correlación, EVMAudit incorpora una
fase adicional orientada a la validación dinámica de determinadas
vulnerabilidades mediante fuzzing.

Para ello, se desarrolló un adaptador específico encargado de
transformar los hallazgos correlacionados en propiedades compatibles con
Echidna, permitiendo complementar los resultados obtenidos mediante
análisis estático y ejecución simbólica.

La incorporación de esta fase responde a uno de los objetivos
específicos planteados en el trabajo, consistente en combinar distintas
técnicas de análisis con el fin de obtener una visión más completa del
estado de seguridad del contrato analizado.

!TODO: referencia cruzada a la Sección 3.2 (objetivos específicos).

\subsubsection{4.7.1. Objetivos del
adaptador}\label{objetivos-del-adaptador}

El módulo desarrollado persigue los siguientes objetivos:

\begin{itemize}
\tightlist
\item
  reutilizar la información obtenida durante la fase de correlación
\item
  generar automáticamente propiedades para Echidna
\item
  reducir la intervención manual del auditor
\item
  proporcionar una fase adicional de validación
\item
  mantener desacopladas las distintas fases del análisis.
\end{itemize}

La existencia de este componente permite integrar el fuzzing dentro del
pipeline general de EVMAudit sin introducir dependencias directas entre
las herramientas utilizadas.

\subsubsection{4.7.2. Generación de
wrappers}\label{generaciuxf3n-de-wrappers}

Echidna requiere que las propiedades de seguridad estén implementadas
como funciones Solidity dentro de un contrato específico.

Con el objetivo de automatizar este proceso, EVMAudit genera
dinámicamente un contrato wrapper que hereda del contrato original e
incorpora las propiedades correspondientes a las vulnerabilidades
detectadas.

El proceso seguido puede resumirse en las siguientes etapas:

\begin{enumerate}
\def\labelenumi{\arabic{enumi}.}
\tightlist
\item
  Recepción de los hallazgos correlacionados.
\item
  Consulta del catálogo SWC.
\item
  Obtención de las plantillas asociadas.
\item
  Sustitución de parámetros.
\item
  Generación del wrapper final.
\item
  Ejecución de Echidna.
\end{enumerate}

Gracias a esta aproximación, la generación de pruebas se realiza de
forma completamente automática.

!TODO: insertar diagrama del proceso de generación de wrappers.

!TODO: insertar referencia a la Figura X.

\subsubsection{4.7.3. Plantillas de
propiedades}\label{plantillas-de-propiedades}

Cada vulnerabilidad soportada dispone de una plantilla específica que
describe la propiedad que deberá ser evaluada por Echidna.

Estas plantillas permiten traducir información abstracta procedente del
análisis estático a estructuras ejecutables compatibles con el motor de
fuzzing.

La utilización de plantillas aporta varias ventajas:

\begin{itemize}
\tightlist
\item
  reutilización de código
\item
  independencia respecto a las herramientas de detección
\item
  facilidad de mantenimiento
\item
  incorporación sencilla de nuevas vulnerabilidades.
\end{itemize}

!TODO: insertar ejemplo de una plantilla real.

\subsubsection{4.7.4. Vulnerabilidades testables y no
testables}\label{vulnerabilidades-testables-y-no-testables}

Durante el desarrollo se observó que no todas las vulnerabilidades
pueden verificarse automáticamente mediante fuzzing.

Por este motivo, el adaptador distingue entre:

\paragraph{Vulnerabilidades testables}\label{vulnerabilidades-testables}

Son aquellas que pueden expresarse mediante propiedades directamente
evaluables por Echidna.

Algunos ejemplos incluyen:

\begin{itemize}
\tightlist
\item
  determinadas restricciones de acceso
\item
  invariantes económicas
\item
  comprobaciones relacionadas con balances
\item
  validaciones de lógica de negocio.
\end{itemize}

\paragraph{Vulnerabilidades no
testables}\label{vulnerabilidades-no-testables}

Corresponden a aquellas situaciones cuya explotación requiere
condiciones externas o escenarios complejos difíciles de reproducir
automáticamente.

Entre ellas destacan:

\begin{itemize}
\tightlist
\item
  ataques de reentrancia
\item
  determinados problemas de interacción entre contratos
\item
  escenarios dependientes del contexto de despliegue.
\end{itemize}

En estos casos, EVMAudit conserva la información asociada a la
vulnerabilidad y genera advertencias específicas, pero no intenta forzar
una validación automática que pudiera producir resultados incorrectos.

Esta decisión se adoptó con el objetivo de evitar conclusiones erróneas
y mantener un comportamiento conservador.

\subsubsection{4.7.5. Integración con
Echidna}\label{integraciuxf3n-con-echidna}

Una vez generado el wrapper correspondiente, el módulo Runner ejecuta
Echidna y recupera los resultados obtenidos durante el proceso de
fuzzing.

La salida generada por la herramienta se procesa posteriormente para
incorporarla al informe final.

De esta manera, EVMAudit consigue integrar tres técnicas de análisis
diferentes:

\begin{itemize}
\tightlist
\item
  análisis estático
\item
  ejecución simbólica
\item
  fuzzing.
\end{itemize}

La combinación de estas aproximaciones permite aprovechar las fortalezas
de cada una de ellas y compensar parcialmente sus limitaciones
individuales.

!TODO: referencia cruzada a la Sección 2.X (análisis estático).

!TODO: referencia cruzada a la Sección 2.X (ejecución simbólica).

!TODO: referencia cruzada a la Sección 2.X (fuzzing).

\subsubsection{4.7.6. Beneficios de la generación
automática}\label{beneficios-de-la-generaciuxf3n-automuxe1tica}

La generación automática de propiedades aporta diversas ventajas:

\begin{itemize}
\tightlist
\item
  reducción del trabajo manual
\item
  mayor reutilización de resultados
\item
  integración transparente con Echidna
\item
  facilidad para incorporar nuevas vulnerabilidades
\item
  mejora del flujo general de análisis.
\end{itemize}

No obstante, la validación automática obtenida no sustituye la revisión
manual realizada por el auditor, sino que constituye una capa adicional
de apoyo.

\section{4.8.Gestion de errores}\label{gestion-de-errores}

\subsection{4.8. Gestión de errores y robustez de la
solución}\label{gestiuxf3n-de-errores-y-robustez-de-la-soluciuxf3n}

Debido a la dependencia de herramientas externas y a la complejidad del
proceso de análisis, durante el diseño de EVMAudit se prestó especial
atención a la gestión de errores y a la robustez del sistema.

La existencia de múltiples componentes externos hace necesario disponer
de mecanismos que permitan detectar y manejar situaciones anómalas sin
comprometer el funcionamiento global de la herramienta.

\subsubsection{4.8.1. Necesidad de una gestión
centralizada}\label{necesidad-de-una-gestiuxf3n-centralizada}

Durante el proceso de análisis pueden producirse diferentes tipos de
errores:

\begin{itemize}
\tightlist
\item
  ausencia de herramientas instaladas
\item
  incompatibilidades entre versiones
\item
  contratos inválidos
\item
  timeouts durante el análisis
\item
  errores internos de las herramientas utilizadas
\item
  problemas durante la generación de wrappers.
\end{itemize}

La gestión individual de cada una de estas situaciones complicaría
considerablemente el mantenimiento del sistema.

Por este motivo se implementó una jerarquía de excepciones propia que
centraliza el tratamiento de los errores.

\subsubsection{4.8.2. Jerarquía de
excepciones}\label{jerarquuxeda-de-excepciones}

La arquitectura implementada incorpora una clase base común a partir de
la cual se derivan los distintos tipos de error utilizados por EVMAudit.

Esta aproximación permite distinguir claramente entre:

\begin{itemize}
\tightlist
\item
  errores de configuración
\item
  errores de análisis
\item
  errores producidos por herramientas externas.
\end{itemize}

La utilización de excepciones específicas simplifica la depuración y
mejora la legibilidad del código.

!TODO: insertar diagrama UML de excepciones.

\subsubsection{4.8.3. Detección de dependencias
externas}\label{detecciuxf3n-de-dependencias-externas}

Antes de ejecutar las herramientas integradas, EVMAudit verifica que las
dependencias necesarias se encuentren correctamente instaladas.

Este mecanismo permite detectar problemas de configuración de forma
temprana y proporcionar mensajes de error más descriptivos.

La validación previa evita que el proceso de análisis falle de forma
inesperada en fases posteriores.

\subsubsection{4.8.4. Gestión de timeouts}\label{gestiuxf3n-de-timeouts}

Algunas herramientas, especialmente aquellas basadas en ejecución
simbólica, pueden presentar tiempos de análisis elevados dependiendo de
la complejidad del contrato.

Con el objetivo de evitar bloqueos indefinidos, EVMAudit incorpora
límites temporales para determinadas operaciones.

La utilización de timeouts permite:

\begin{itemize}
\tightlist
\item
  mejorar la estabilidad del sistema
\item
  evitar consumos excesivos de recursos
\item
  garantizar la finalización del análisis.
\end{itemize}

\subsubsection{4.8.5. Preservación de resultados ante
errores}\label{preservaciuxf3n-de-resultados-ante-errores}

Cuando se produce un fallo durante alguna fase del análisis, EVMAudit
intenta conservar la información obtenida hasta ese momento.

Esta estrategia permite:

\begin{itemize}
\tightlist
\item
  facilitar la depuración
\item
  conservar evidencias parciales
\item
  evitar la pérdida completa del trabajo realizado.
\end{itemize}

La preservación de resultados resulta especialmente relevante en
entornos de auditoría, donde incluso los análisis incompletos pueden
aportar información útil.

\subsubsection{4.8.6. Robustez frente a comportamientos específicos de
las
herramientas}\label{robustez-frente-a-comportamientos-especuxedficos-de-las-herramientas}

Durante el desarrollo se detectaron determinados comportamientos
particulares en las herramientas integradas.

Entre ellos destacan:

\begin{itemize}
\tightlist
\item
  códigos de retorno no convencionales
\item
  formatos de salida inconsistentes
\item
  diferencias entre versiones
\item
  mensajes mezclados con datos estructurados.
\end{itemize}

Para hacer frente a estas situaciones se implementaron mecanismos
específicos de adaptación y validación.

Estas medidas permitieron integrar las herramientas seleccionadas sin
modificar su funcionamiento interno.

\subsubsection{4.8.7. Beneficios de la estrategia
adoptada}\label{beneficios-de-la-estrategia-adoptada}

La incorporación de mecanismos de gestión de errores proporciona
diversas ventajas:

\begin{itemize}
\tightlist
\item
  mayor estabilidad
\item
  mejora de la experiencia de usuario
\item
  facilidad de mantenimiento
\item
  simplificación de la depuración
\item
  incremento de la robustez general del sistema.
\end{itemize}

La existencia de una capa de gestión de errores independiente contribuye
además a mantener desacoplados los distintos módulos que forman la
arquitectura de EVMAudit.

!TODO: insertar referencia cruzada al módulo Runner.

!TODO: insertar referencia cruzada al diagrama de arquitectura general.

\section{4.9.Integración y
distribución}\label{integraciuxf3n-y-distribuciuxf3n}

\subsection{4.9. Distribución e integración de la
solución}\label{distribuciuxf3n-e-integraciuxf3n-de-la-soluciuxf3n}

Uno de los objetivos perseguidos durante el desarrollo de EVMAudit fue
diseñar una solución desacoplada de cualquier interfaz concreta y
suficientemente flexible para poder ser reutilizada en distintos
contextos.

Por este motivo, la lógica principal de análisis se implementó como una
librería independiente, separando completamente las funcionalidades
relacionadas con la detección y procesamiento de vulnerabilidades de los
mecanismos de interacción con el usuario.

Esta decisión permitió demostrar la reutilización de la solución
desarrollada mediante su integración en otros componentes software sin
necesidad de modificar el núcleo de la herramienta.

\subsubsection{4.9.1. Distribución de la
librería}\label{distribuciuxf3n-de-la-libreruxeda}

La implementación de EVMAudit como una librería independiente permite
desacoplar el motor de análisis de cualquier entorno concreto de
ejecución.

Este enfoque proporciona diversas ventajas:

\begin{itemize}
\tightlist
\item
  reutilización de la herramienta desde otros proyectos
\item
  separación entre lógica de negocio e interfaces de usuario
\item
  facilidad de mantenimiento
\item
  posibilidad de incorporar nuevas interfaces en el futuro.
\end{itemize}

La distribución de la librería permite que el proceso de análisis pueda
ejecutarse desde aplicaciones externas sin necesidad de duplicar
funcionalidades.

!TODO: referencia cruzada al Anexo A (publicación en PyPI).

\subsubsection{4.9.2. Aplicación web como caso de
uso}\label{aplicaciuxf3n-web-como-caso-de-uso}

Con el objetivo de validar la reutilización de la arquitectura
desarrollada, se implementó una aplicación web que utiliza EVMAudit como
motor de análisis.

La aplicación no implementa mecanismos propios de detección de
vulnerabilidades, sino que delega completamente las tareas de análisis
en la librería desarrollada.

Esta aproximación permite mantener una única implementación del proceso
de análisis y garantiza la consistencia de los resultados obtenidos.

La aplicación actúa, por tanto, como una capa de presentación encargada
de:

\begin{itemize}
\tightlist
\item
  recibir contratos inteligentes proporcionados por el usuario
\item
  invocar las funciones de EVMAudit
\item
  mostrar el progreso del análisis
\item
  presentar los resultados obtenidos.
\end{itemize}

De esta forma, la aplicación web constituye una demostración práctica de
la reutilización de la solución desarrollada.

\subsubsection{4.9.3. Flujo de funcionamiento de la
aplicación}\label{flujo-de-funcionamiento-de-la-aplicaciuxf3n}

La aplicación desarrollada permite dos mecanismos de entrada:

\begin{itemize}
\tightlist
\item
  carga de contratos Solidity mediante fichero
\item
  introducción directa del código fuente a través de la interfaz web.
\end{itemize}

Una vez recibido el contrato, la aplicación inicia el proceso de
análisis y ejecuta internamente el pipeline implementado por EVMAudit.

Las distintas fases del proceso se ejecutan de forma secuencial:

\begin{enumerate}
\def\labelenumi{\arabic{enumi}.}
\tightlist
\item
  ejecución de Slither
\item
  ejecución de Mythril
\item
  normalización de resultados
\item
  correlación de vulnerabilidades
\item
  generación de propiedades
\item
  ejecución de Echidna
\item
  generación del informe final.
\end{enumerate}

Durante la ejecución, la aplicación informa al usuario del progreso
alcanzado y permite recuperar el informe una vez finalizado el análisis.

Este comportamiento demuestra que la arquitectura modular adoptada
permite integrar EVMAudit en aplicaciones externas sin modificar la
lógica interna del sistema.

!TODO: insertar diagrama simplificado de la aplicación web.

!TODO: insertar captura de pantalla de la interfaz.

!TODO: referencia cruzada al Anexo B (aplicación web).

\subsubsection{4.9.4. Separación entre presentación y lógica de
análisis}\label{separaciuxf3n-entre-presentaciuxf3n-y-luxf3gica-de-anuxe1lisis}

La decisión de implementar EVMAudit como una librería independiente
permite mantener separadas las responsabilidades del sistema.

Mientras que la librería se encarga de:

\begin{itemize}
\tightlist
\item
  ejecutar herramientas externas
\item
  procesar resultados
\item
  correlacionar vulnerabilidades
\item
  generar informes
\end{itemize}

la aplicación web únicamente proporciona los mecanismos de interacción
con el usuario.

Esta separación favorece:

\begin{itemize}
\tightlist
\item
  la mantenibilidad del sistema
\item
  la reutilización de la solución
\item
  la incorporación de nuevas interfaces
\item
  la evolución independiente de cada componente.
\end{itemize}

En consecuencia, la aplicación web debe entenderse como una demostración
de la capacidad de integración de EVMAudit y no como la contribución
principal del trabajo.

\subsubsection{4.9.5. Consideraciones sobre
infraestructura}\label{consideraciones-sobre-infraestructura}

Los aspectos relacionados con la publicación de la librería, el
despliegue de la aplicación y la infraestructura utilizada no
constituyen el núcleo de la contribución desarrollada en este Trabajo
Fin de Máster.

Por este motivo, dichos elementos se incluyen en los anexos del
documento.

!TODO: referencia cruzada al Anexo A (PyPI).

!TODO: referencia cruzada al Anexo B (Aplicación web).

!TODO: referencia cruzada al Anexo C (Docker).

!TODO: referencia cruzada al Anexo D (Despliegue).

\section{4.10. Flujo completo}\label{flujo-completo}

\subsection{4.10. Flujo completo del
análisis}\label{flujo-completo-del-anuxe1lisis}

Una vez descritos los distintos módulos que componen EVMAudit, resulta
posible analizar el funcionamiento global de la solución.

El proceso completo implementado por la herramienta se compone de varias
fases consecutivas que permiten transformar un contrato inteligente en
un informe consolidado de vulnerabilidades.

\subsubsection{4.10.1. Recepción del
contrato}\label{recepciuxf3n-del-contrato}

El proceso comienza con la recepción del contrato inteligente que se
desea analizar.

El contrato puede proceder de distintos orígenes:

\begin{itemize}
\tightlist
\item
  ejecución directa de la librería
\item
  aplicación web desarrollada
\item
  futuras integraciones externas.
\end{itemize}

Tras recibir el contrato, EVMAudit detecta automáticamente el nombre del
contrato y configura la versión adecuada del compilador Solidity.

\subsubsection{4.10.2. Obtención de resultados mediante Slither y
Mythril}\label{obtenciuxf3n-de-resultados-mediante-slither-y-mythril}

La siguiente fase consiste en ejecutar las herramientas de análisis
principales.

En primer lugar se ejecuta Slither, aprovechando las ventajas del
análisis estático.

Posteriormente se ejecuta Mythril, incorporando las capacidades
derivadas de la ejecución simbólica.

La combinación de ambas técnicas permite obtener una cobertura superior
a la proporcionada por cada herramienta de forma individual.

!TODO: referencia cruzada a la Sección 2.X (herramientas de análisis).

\subsubsection{4.10.3. Normalización de
resultados}\label{normalizaciuxf3n-de-resultados}

Las salidas generadas por las herramientas se transforman posteriormente
a una estructura común.

Este proceso elimina las diferencias existentes entre ambas herramientas
y permite trabajar con una representación homogénea de las
vulnerabilidades.

La normalización constituye el paso previo necesario para aplicar el
mecanismo de correlación.

\subsubsection{4.10.4. Correlación de
vulnerabilidades}\label{correlaciuxf3n-de-vulnerabilidades}

Una vez normalizados los resultados, EVMAudit agrupa aquellas
vulnerabilidades equivalentes detectadas por varias herramientas.

La correlación permite:

\begin{itemize}
\tightlist
\item
  reducir duplicidades
\item
  aumentar la confianza de determinados hallazgos
\item
  simplificar la interpretación del informe final.
\end{itemize}

Esta fase representa uno de los principales elementos diferenciadores de
la solución desarrollada.

\subsubsection{4.10.5. Generación de propiedades y
fuzzing}\label{generaciuxf3n-de-propiedades-y-fuzzing}

A partir de las vulnerabilidades correlacionadas, el sistema consulta el
catálogo SWC y genera automáticamente las propiedades necesarias para
ejecutar Echidna.

Posteriormente se construye un wrapper Solidity específico y se realiza
una fase adicional de validación mediante fuzzing.

Este proceso proporciona información complementaria sobre determinadas
vulnerabilidades detectadas durante las fases anteriores.

\subsubsection{4.10.6. Generación del informe
final}\label{generaciuxf3n-del-informe-final}

Finalmente, toda la información obtenida durante el análisis es
consolidada en un informe único.

El informe incluye:

\begin{itemize}
\tightlist
\item
  vulnerabilidades detectadas
\item
  herramientas que las han identificado
\item
  nivel de confianza asociado
\item
  resultados del fuzzing
\item
  metadatos adicionales.
\end{itemize}

De este modo, el auditor dispone de una visión unificada del estado de
seguridad del contrato analizado.

\subsubsection{4.10.7. Resumen del proceso}\label{resumen-del-proceso}

El flujo general implementado por EVMAudit puede resumirse en las
siguientes etapas:

\begin{enumerate}
\def\labelenumi{\arabic{enumi}.}
\tightlist
\item
  Recepción del contrato.
\item
  Configuración del entorno.
\item
  Ejecución de Slither.
\item
  Ejecución de Mythril.
\item
  Normalización.
\item
  Correlación.
\item
  Consulta del catálogo SWC.
\item
  Generación del wrapper.
\item
  Ejecución de Echidna.
\item
  Generación del informe.
\end{enumerate}

La secuencia completa permite combinar distintas técnicas de análisis
dentro de una arquitectura unificada y automatizada.

!TODO: insertar diagrama de secuencia del pipeline completo.

!TODO: insertar diagrama de flujo del proceso.

!TODO: insertar referencia a la Figura X.

!TODO: insertar referencia cruzada a las secciones 4.3 a 4.9.

\section{4.11.Evaluación experimental}\label{evaluaciuxf3n-experimental}

\subsection{4.11. Flujo completo del
análisis}\label{flujo-completo-del-anuxe1lisis}

Una vez descritos los distintos módulos que componen EVMAudit, resulta
posible analizar el funcionamiento global de la solución.

El proceso completo implementado por la herramienta se compone de varias
fases consecutivas que permiten transformar un contrato inteligente en
un informe consolidado de vulnerabilidades.

\subsubsection{4.11.1. Recepción del
contrato}\label{recepciuxf3n-del-contrato}

El proceso comienza con la recepción del contrato inteligente que se
desea analizar.

El contrato puede proceder de distintos orígenes:

\begin{itemize}
\tightlist
\item
  ejecución directa de la librería
\item
  aplicación web desarrollada
\item
  futuras integraciones externas.
\end{itemize}

Tras recibir el contrato, EVMAudit detecta automáticamente el nombre del
contrato y configura la versión adecuada del compilador Solidity.

\subsubsection{4.11.2. Obtención de resultados mediante Slither y
Mythril}\label{obtenciuxf3n-de-resultados-mediante-slither-y-mythril}

La siguiente fase consiste en ejecutar las herramientas de análisis
principales.

En primer lugar se ejecuta Slither, aprovechando las ventajas del
análisis estático.

Posteriormente se ejecuta Mythril, incorporando las capacidades
derivadas de la ejecución simbólica.

La combinación de ambas técnicas permite obtener una cobertura superior
a la proporcionada por cada herramienta de forma individual.

!TODO: referencia cruzada a la Sección 2.X (herramientas de análisis).

\subsubsection{4.11.3. Normalización de
resultados}\label{normalizaciuxf3n-de-resultados}

Las salidas generadas por las herramientas se transforman posteriormente
a una estructura común.

Este proceso elimina las diferencias existentes entre ambas herramientas
y permite trabajar con una representación homogénea de las
vulnerabilidades.

La normalización constituye el paso previo necesario para aplicar el
mecanismo de correlación.

\subsubsection{4.11.4. Correlación de
vulnerabilidades}\label{correlaciuxf3n-de-vulnerabilidades}

Una vez normalizados los resultados, EVMAudit agrupa aquellas
vulnerabilidades equivalentes detectadas por varias herramientas.

La correlación permite:

\begin{itemize}
\tightlist
\item
  reducir duplicidades
\item
  aumentar la confianza de determinados hallazgos
\item
  simplificar la interpretación del informe final.
\end{itemize}

Esta fase representa uno de los principales elementos diferenciadores de
la solución desarrollada.

\subsubsection{4.11.5. Generación de propiedades y
fuzzing}\label{generaciuxf3n-de-propiedades-y-fuzzing}

A partir de las vulnerabilidades correlacionadas, el sistema consulta el
catálogo SWC y genera automáticamente las propiedades necesarias para
ejecutar Echidna.

Posteriormente se construye un wrapper Solidity específico y se realiza
una fase adicional de validación mediante fuzzing.

Este proceso proporciona información complementaria sobre determinadas
vulnerabilidades detectadas durante las fases anteriores.

\subsubsection{4.11.6. Generación del informe
final}\label{generaciuxf3n-del-informe-final}

Finalmente, toda la información obtenida durante el análisis es
consolidada en un informe único.

El informe incluye:

\begin{itemize}
\tightlist
\item
  vulnerabilidades detectadas
\item
  herramientas que las han identificado
\item
  nivel de confianza asociado
\item
  resultados del fuzzing
\item
  metadatos adicionales.
\end{itemize}

De este modo, el auditor dispone de una visión unificada del estado de
seguridad del contrato analizado.

\subsubsection{4.11.7. Resumen del proceso}\label{resumen-del-proceso}

El flujo general implementado por EVMAudit puede resumirse en las
siguientes etapas:

\begin{enumerate}
\def\labelenumi{\arabic{enumi}.}
\tightlist
\item
  Recepción del contrato.
\item
  Configuración del entorno.
\item
  Ejecución de Slither.
\item
  Ejecución de Mythril.
\item
  Normalización.
\item
  Correlación.
\item
  Consulta del catálogo SWC.
\item
  Generación del wrapper.
\item
  Ejecución de Echidna.
\item
  Generación del informe.
\end{enumerate}

La secuencia completa permite combinar distintas técnicas de análisis
dentro de una arquitectura unificada y automatizada.

!TODO: insertar diagrama de secuencia del pipeline completo.

!TODO: insertar diagrama de flujo del proceso.

!TODO: insertar referencia a la Figura X.

!TODO: insertar referencia cruzada a las secciones 4.3 a 4.9.

\subsubsection{4.11.X. Evaluación del contrato con una vulnerabilidad
conocida: {[}NOMBRE DEL
CONTRATO{]}}\label{x.-evaluaciuxf3n-del-contrato-con-una-vulnerabilidad-conocida-nombre-del-contrato}

\paragraph{Descripción del contrato}\label{descripciuxf3n-del-contrato}

El contrato \textbf{{[}NOMBRE DEL CONTRATO{]}} fue seleccionado debido a
que presenta una vulnerabilidad conocida relacionada con \textbf{{[}TIPO
DE VULNERABILIDAD{]}}.

Este contrato pertenece a \textbf{{[}FUENTE DEL CONTRATO{]}} y fue
incluido en la evaluación con el objetivo de analizar el comportamiento
de EVMAudit frente a escenarios representativos.

La vulnerabilidad esperada en este caso corresponde a:

\begin{itemize}
\tightlist
\item
  TODO.
\end{itemize}

!TODO: referencia cruzada a la Sección 2.X (descripción de la
vulnerabilidad).

\begin{center}\rule{0.5\linewidth}{0.5pt}\end{center}

\paragraph{Resultados obtenidos por las
herramientas}\label{resultados-obtenidos-por-las-herramientas}

En primer lugar, el contrato fue procesado individualmente por Slither y
Mythril.

Los resultados obtenidos se resumen en la Tabla X.

{\def\LTcaptype{none} % do not increment counter
\begin{longtable}[]{@{}ll@{}}
\toprule\noalign{}
Herramienta & Vulnerabilidades detectadas \\
\midrule\noalign{}
\endhead
\bottomrule\noalign{}
\endlastfoot
Slither & TODO \\
Mythril & TODO \\
\end{longtable}
}

!TODO: insertar referencia a la Tabla X.

\begin{center}\rule{0.5\linewidth}{0.5pt}\end{center}

\paragraph{Resultados tras la
correlación}\label{resultados-tras-la-correlaciuxf3n}

Una vez aplicados los mecanismos de normalización y correlación
implementados por EVMAudit, se obtuvo un total de \textbf{TODO}
vulnerabilidades consolidadas.

La Tabla X muestra la relación entre los hallazgos originales y los
resultados finales.

{\def\LTcaptype{none} % do not increment counter
\begin{longtable}[]{@{}ll@{}}
\toprule\noalign{}
Métrica & Valor \\
\midrule\noalign{}
\endhead
\bottomrule\noalign{}
\endlastfoot
Hallazgos originales & TODO \\
Hallazgos correlacionados & TODO \\
Vulnerabilidades confirmadas & TODO \\
Reducción de duplicidades & TODO \\
\end{longtable}
}

!TODO: insertar referencia a la Tabla X.

\begin{center}\rule{0.5\linewidth}{0.5pt}\end{center}

\paragraph{Resultados del fuzzing mediante
Echidna}\label{resultados-del-fuzzing-mediante-echidna}

A partir de las vulnerabilidades obtenidas se generaron automáticamente
las propiedades correspondientes para Echidna.

Los resultados obtenidos fueron los siguientes:

\begin{itemize}
\tightlist
\item
  Propiedades generadas: TODO.
\item
  Propiedades ejecutadas correctamente: TODO.
\item
  Vulnerabilidades verificadas: TODO.
\item
  Limitaciones observadas: TODO.
\end{itemize}

!TODO: insertar captura o fragmento representativo del resultado de
Echidna.

\begin{center}\rule{0.5\linewidth}{0.5pt}\end{center}

\paragraph{Discusión}\label{discusiuxf3n}

Los resultados obtenidos muestran que \textbf{{[}OBSERVACIÓN
PRINCIPAL{]}}.

En particular:

\begin{itemize}
\tightlist
\item
  TODO.
\item
  TODO.
\item
  TODO.
\end{itemize}

Este caso pone de manifiesto \textbf{{[}CONCLUSIÓN PRINCIPAL DEL
CONTRATO{]}}.

\subsubsection{4.11.X. Evaluación del contrato con varias
vulnerabilidades: {[}NOMBRE DEL
CONTRATO{]}}\label{x.-evaluaciuxf3n-del-contrato-con-varias-vulnerabilidades-nombre-del-contrato}

\paragraph{Descripción del
contrato}\label{descripciuxf3n-del-contrato-1}

El contrato \textbf{{[}NOMBRE{]}} incorpora una vulnerabilidad
relacionada con \textbf{{[}TIPO{]}}, lo que permite estudiar la
capacidad de detección de las herramientas integradas cuando existen
varias debilidades simultáneas.

!TODO: describir brevemente el funcionamiento del contrato.

\begin{center}\rule{0.5\linewidth}{0.5pt}\end{center}

\paragraph{Comparativa entre
herramientas}\label{comparativa-entre-herramientas}

La Tabla X resume las vulnerabilidades detectadas por cada herramienta.

{\def\LTcaptype{none} % do not increment counter
\begin{longtable}[]{@{}lcc@{}}
\toprule\noalign{}
Vulnerabilidad & Slither & Mythril \\
\midrule\noalign{}
\endhead
\bottomrule\noalign{}
\endlastfoot
TODO & {[}X{]} & {[}X{]} \\
TODO & {[}X{]} & {[}-{]} \\
TODO & {[}-{]} & {[}X{]} \\
\end{longtable}
}

Los resultados muestran que ambas herramientas presentan comportamientos
complementarios.

\begin{center}\rule{0.5\linewidth}{0.5pt}\end{center}

\paragraph{Efecto del mecanismo de
correlación}\label{efecto-del-mecanismo-de-correlaciuxf3n}

La aplicación del proceso de correlación permitió:

\begin{itemize}
\tightlist
\item
  eliminar duplicidades;
\item
  unificar nomenclaturas;
\item
  aumentar la confianza de determinados hallazgos.
\end{itemize}

!TODO: insertar número real de vulnerabilidades correlacionadas.

\begin{center}\rule{0.5\linewidth}{0.5pt}\end{center}

\paragraph{Resultados del fuzzing}\label{resultados-del-fuzzing}

El proceso de generación automática produjo un total de \textbf{TODO}
propiedades.

Durante la ejecución de Echidna se observó:

\begin{itemize}
\tightlist
\item
  TODO.
\end{itemize}

!TODO: describir vulnerabilidades no testables.

\begin{center}\rule{0.5\linewidth}{0.5pt}\end{center}

\paragraph{Discusión}\label{discusiuxf3n-1}

Este contrato evidencia que las distintas técnicas de análisis empleadas
por EVMAudit permiten obtener información complementaria y mejorar la
cobertura respecto al uso individual de cada herramienta.

\subsubsection{4.11.X. Evaluación del contrato real {[}NOMBRE DEL
CONTRATO{]}}\label{x.-evaluaciuxf3n-del-contrato-real-nombre-del-contrato}

\paragraph{Descripción del
contrato}\label{descripciuxf3n-del-contrato-2}

Además de los contratos vulnerables de referencia, se incluyó el
contrato \textbf{{[}NOMBRE{]}}, obtenido de \textbf{{[}REPOSITORIO O
FUENTE{]}}, con el objetivo de analizar el comportamiento de EVMAudit en
un escenario más próximo a un entorno real.

!TODO: describir brevemente la finalidad del contrato.

\begin{center}\rule{0.5\linewidth}{0.5pt}\end{center}

\paragraph{Resultados obtenidos}\label{resultados-obtenidos}

La Tabla X resume las vulnerabilidades detectadas durante el análisis.

{\def\LTcaptype{none} % do not increment counter
\begin{longtable}[]{@{}ll@{}}
\toprule\noalign{}
Métrica & Valor \\
\midrule\noalign{}
\endhead
\bottomrule\noalign{}
\endlastfoot
Vulnerabilidades detectadas por Slither & TODO \\
Vulnerabilidades detectadas por Mythril & TODO \\
Vulnerabilidades finales & TODO \\
Vulnerabilidades confirmadas & TODO \\
\end{longtable}
}

\begin{center}\rule{0.5\linewidth}{0.5pt}\end{center}

\paragraph{Observaciones sobre la
correlación}\label{observaciones-sobre-la-correlaciuxf3n}

En este caso se observó que:

\begin{itemize}
\tightlist
\item
  TODO.
\end{itemize}

Asimismo, se identificaron diferencias entre ambas herramientas en
relación con:

\begin{itemize}
\tightlist
\item
  TODO.
\end{itemize}

\begin{center}\rule{0.5\linewidth}{0.5pt}\end{center}

\paragraph{Resultados del fuzzing}\label{resultados-del-fuzzing-1}

La fase de fuzzing permitió:

\begin{itemize}
\tightlist
\item
  TODO.
\end{itemize}

No obstante, algunas vulnerabilidades no pudieron validarse
automáticamente debido a:

\begin{itemize}
\tightlist
\item
  TODO.
\end{itemize}

\begin{center}\rule{0.5\linewidth}{0.5pt}\end{center}

\paragraph{Discusión}\label{discusiuxf3n-2}

El análisis realizado demuestra que EVMAudit puede aplicarse también
sobre contratos reales y no únicamente sobre ejemplos académicos
diseñados específicamente para contener vulnerabilidades.

!TODO: indicar limitaciones observadas durante el análisis.


\chapter{Conclusiones y trabajo futuro}
\section{5. Conclusiones y trabajo
futuro}\label{conclusiones-y-trabajo-futuro}

\subsection{5.1 Conclusiones}\label{conclusiones}

\subsection{5.2 Limitaciones de la
solución}\label{limitaciones-de-la-soluciuxf3n}

\begin{itemize}
\item
  \textbf{Dependencia de herramientas externas.} EVMAudit depende del
  correcto funcionamiento de Slither, Mythril y Echidna. Cualquier
  modificación en dichas herramientas puede afectar al comportamiento
  del sistema.
\item
  \textbf{Correlación basada en reglas.} La correlación implementada
  utiliza únicamente:

  \begin{itemize}
  \tightlist
  \item
    contrato
  \item
    función
  \item
    SWC
  \end{itemize}

  Esto puede provocar que determinadas vulnerabilidades relacionadas no
  se agrupen correctamente.
\item
  \textbf{Falsos positivos y falsos negativos.} La herramienta no
  elimina completamente los falsos positivos inherentes a las
  herramientas integradas.
\item
  \textbf{Limitaciones del fuzzing.} Determinadas vulnerabilidades, como
  la reentrancia, requieren contratos atacantes específicos y no pueden
  validarse automáticamente.
\item
  \textbf{Ausencia de análisis inter-contrato.} Actualmente la
  herramienta analiza cada contrato de forma independiente.
\item
  \textbf{Priorización limitada.} La priorización se basa en la
  severidad y en el número de herramientas que detectan la
  vulnerabilidad, sin utilizar métricas más avanzadas como CVSS.
\end{itemize}

\subsection{5.3 Líneas de trabajo
futuro}\label{luxedneas-de-trabajo-futuro}

\begin{itemize}
\tightlist
\item
  incorporar nuevas herramientas
\item
  análisis inter-contrato\\
\item
  CVSS
\item
  Machine Learning
\item
  Semgrep
\item
  Manticore
\item
  dashboard web
\end{itemize}


\bibliographystyle{apacite}
\bibliography{bibliografia}

\appendix

\chapter{Ejemplos de vulnerabilidades}
\section{ANEXO A.~Ejemplos de vulnerabilidades en contratos
inteligentes}\label{anexo-a.-ejemplos-de-vulnerabilidades-en-contratos-inteligentes}

Este anexo presenta ejemplos simplificados de vulnerabilidades
representativas en contratos inteligentes desarrollados en Solidity. El
objetivo es ilustrar de forma práctica los principales tipos de
debilidades descritas en la sección 4.4, facilitando su comprensión y su
posterior detección mediante herramientas automáticas.

Cada ejemplo incluye: descripción, fragmento de código vulnerable,
impacto y estrategia de mitigación.

\subsection{ANEXO A.1. Vulnerabilidades técnicas de
ejecución}\label{anexo-a.1.-vulnerabilidades-tuxe9cnicas-de-ejecuciuxf3n}

\subsubsection{ANEXO A.1.1. Reentrancy}\label{anexo-a.1.1.-reentrancy}

La vulnerabilidad de
\href{https://www.cyfrin.io/blog/what-is-a-reentrancy-attack-solidity-smart-contracts}{\textbf{reentrancy}}
se produce cuando un contrato realiza una llamada externa antes de
actualizar su estado interno, permitiendo que el contrato receptor
reingrese en la función original en un estado inconsistente.

Código vulnerable

\begin{Shaded}
\begin{Highlighting}[]
\NormalTok{contract VulnerableBank \{}

    \FunctionTok{mapping}\NormalTok{(address }\KeywordTok{=\textgreater{}}\NormalTok{ uint256) }\KeywordTok{public}\NormalTok{ balances}\OperatorTok{;}

    \KeywordTok{function} \FunctionTok{deposit}\NormalTok{() }\KeywordTok{public}\NormalTok{ payable \{}

\NormalTok{        balances[msg}\OperatorTok{.}\AttributeTok{sender}\NormalTok{] }\OperatorTok{+=}\NormalTok{ msg}\OperatorTok{.}\AttributeTok{value}\OperatorTok{;}

\NormalTok{    \}}

    \KeywordTok{function} \FunctionTok{withdraw}\NormalTok{(uint256 amount) }\KeywordTok{public}\NormalTok{ \{}

        \PreprocessorTok{require}\NormalTok{(balances[msg}\OperatorTok{.}\AttributeTok{sender}\NormalTok{] }\OperatorTok{\textgreater{}=}\NormalTok{ amount)}\OperatorTok{;}

\FunctionTok{       }\NormalTok{ (bool success}\OperatorTok{,}\NormalTok{ ) }\OperatorTok{=}\NormalTok{ msg}\OperatorTok{.}\AttributeTok{sender}\OperatorTok{.}\AttributeTok{call}\NormalTok{\{}\DataTypeTok{value}\OperatorTok{:}\NormalTok{ amount\}(}\StringTok{""}\NormalTok{)}\OperatorTok{;}

        \PreprocessorTok{require}\NormalTok{(success)}\OperatorTok{;}

\NormalTok{        balances[msg}\OperatorTok{.}\AttributeTok{sender}\NormalTok{] }\OperatorTok{{-}=}\NormalTok{ amount}\OperatorTok{;}

\NormalTok{    \}}

\NormalTok{\}}
\end{Highlighting}
\end{Shaded}

\begin{itemize}
\tightlist
\item
  \textbf{Impacto}: Un atacante puede drenar fondos repitiendo la
  llamada antes de que el balance sea actualizado.
\item
  Mitigación:

  \begin{itemize}
  \tightlist
  \item
    Patrón
    \href{https://fravoll.github.io/solidity-patterns/checks_effects_interactions.html}{\emph{Checks-Effects-Interactions}}
  \item
    Uso de ReentrancyGuard
  \item
    Actualizar el estado antes de llamadas externas
  \end{itemize}
\end{itemize}

\subsubsection{ANEXO A.1.2. Integer Overflow /
Underflow}\label{anexo-a.1.2.-integer-overflow-underflow}

Errores aritméticos que provocan desbordamientos en operaciones enteras.
Aunque mitigados en Solidity \textgreater= 0.8, siguen siendo relevantes
en bloques unchecked.

Código vulnerable

\begin{Shaded}
\begin{Highlighting}[]
\KeywordTok{function} \FunctionTok{increment}\NormalTok{(uint256 x) }\KeywordTok{public}\NormalTok{ pure }\FunctionTok{returns}\NormalTok{ (uint256) \{}

\NormalTok{    unchecked \{}

        \ControlFlowTok{return}\NormalTok{ x }\OperatorTok{+} \DecValTok{1}\OperatorTok{;}

\NormalTok{    \}}

\NormalTok{\}}
\end{Highlighting}
\end{Shaded}

\begin{itemize}
\tightlist
\item
  \textbf{Impacto}: Puede alterar balances o condiciones lógicas
  críticas.
\item
  Mitigación:

  \begin{itemize}
  \tightlist
  \item
    Evitar unchecked salvo casos justificados
  \item
    Uso de validaciones explícitas
  \end{itemize}
\end{itemize}

\subsubsection{ANEXO A.1.3. Uso inseguro de
delegatecall}\label{anexo-a.1.3.-uso-inseguro-de-delegatecall}

La función delegatecall ejecuta código externo en el contexto de
almacenamiento del contrato llamador.

Código vulnerable

\begin{Shaded}
\begin{Highlighting}[]
\NormalTok{contract }\BuiltInTok{Proxy}\NormalTok{ \{}

\NormalTok{    address }\KeywordTok{public}\NormalTok{ implementation}\OperatorTok{;}

    \KeywordTok{function} \FunctionTok{execute}\NormalTok{(bytes memory data) }\KeywordTok{public}\NormalTok{ \{}

\FunctionTok{       }\NormalTok{ (bool success}\OperatorTok{,}\NormalTok{ ) }\OperatorTok{=}\NormalTok{ implementation}\OperatorTok{.}\FunctionTok{delegatecall}\NormalTok{(data)}\OperatorTok{;}

        \PreprocessorTok{require}\NormalTok{(success)}\OperatorTok{;}

\NormalTok{    \}}

\NormalTok{\}}
\end{Highlighting}
\end{Shaded}

\begin{itemize}
\tightlist
\item
  \textbf{Impacto}: Compromiso total del almacenamiento del contrato.
\item
  Mitigación:

  \begin{itemize}
  \tightlist
  \item
    Control estricto de la dirección implementation
  \item
    Uso de patrones proxy auditados
    (\href{https://eips.ethereum.org/EIPS/eip-1967}{EIP-1967},
    \href{https://docs.openzeppelin.com/contracts-stylus/uups-proxy}{UUPS})
  \end{itemize}
\end{itemize}

\subsubsection{ANEXO A.1.4. Denegación de servicio
(DoS)}\label{anexo-a.1.4.-denegaciuxf3n-de-servicio-dos}

Bloqueo de ejecución debido a fallos en llamadas externas o estructuras
no acotadas.

Código vulnerable

\begin{Shaded}
\begin{Highlighting}[]
\KeywordTok{function} \FunctionTok{payout}\NormalTok{(address[] memory recipients) }\KeywordTok{public}\NormalTok{ \{}

    \ControlFlowTok{for}\NormalTok{ (uint i }\OperatorTok{=} \DecValTok{0}\OperatorTok{;}\NormalTok{ i }\OperatorTok{\textless{}}\NormalTok{ recipients}\OperatorTok{.}\AttributeTok{length}\OperatorTok{;}\NormalTok{ i}\OperatorTok{++}\NormalTok{) \{}

        \FunctionTok{payable}\NormalTok{(recipients[i])}\OperatorTok{.}\FunctionTok{transfer}\NormalTok{(}\DecValTok{1}\NormalTok{ ether)}\OperatorTok{;}

\NormalTok{    \}}

\NormalTok{\}}
\end{Highlighting}
\end{Shaded}

\begin{itemize}
\tightlist
\item
  \textbf{Impacto}: Un solo fallo revierte toda la operación.
\item
  Mitigación

  \begin{itemize}
  \tightlist
  \item
    Uso de
    \href{https://medium.com/@markojauregui/the-pull-over-push-model-in-solidity-a-secure-pattern-for-fund-withdrawals-10c2e6628626}{patrón
    \emph{pull over push}}
  \item
    Evitar bucles dependientes de input externo
  \end{itemize}
\end{itemize}

\subsection{ANEXO A.2. Vulnerabilidades técnicas de control y
privilegios}\label{anexo-a.2.-vulnerabilidades-tuxe9cnicas-de-control-y-privilegios}

\subsubsection{ANEXO A.2.1. Falta de control de
acceso}\label{anexo-a.2.1.-falta-de-control-de-acceso}

Funciones críticas accesibles por cualquier usuario.

Código vulnerable

\begin{Shaded}
\begin{Highlighting}[]
\NormalTok{contract Ownable \{}

\NormalTok{    address }\KeywordTok{public}\NormalTok{ owner}\OperatorTok{;}

    \KeywordTok{function} \FunctionTok{withdrawAll}\NormalTok{() }\KeywordTok{public}\NormalTok{ \{}

        \FunctionTok{payable}\NormalTok{(msg}\OperatorTok{.}\AttributeTok{sender}\NormalTok{)}\OperatorTok{.}\FunctionTok{transfer}\NormalTok{(}\FunctionTok{address}\NormalTok{(}\KeywordTok{this}\NormalTok{)}\OperatorTok{.}\AttributeTok{balance}\NormalTok{)}\OperatorTok{;}

\NormalTok{    \}}

\NormalTok{\}}
\end{Highlighting}
\end{Shaded}

\begin{itemize}
\tightlist
\item
  \textbf{Impacto}: Pérdida total de fondos.
\item
  Mitigación:

  \begin{itemize}
  \tightlist
  \item
    Uso de onlyOwner
  \item
    Librerías como
    \href{https://docs.openzeppelin.com/contracts/5.x/access-control}{OpenZeppelin
    AccessControl}
  \end{itemize}
\end{itemize}

\subsubsection{ANEXO A.2.2. Uso de
tx.origin}\label{anexo-a.2.2.-uso-de-tx.origin}

Uso incorrecto de tx.origin para autenticación.

Código vulnerable

\begin{Shaded}
\begin{Highlighting}[]
\KeywordTok{function} \FunctionTok{withdraw}\NormalTok{() }\KeywordTok{public}\NormalTok{ \{}

    \PreprocessorTok{require}\NormalTok{(tx}\OperatorTok{.}\AttributeTok{origin} \OperatorTok{==}\NormalTok{ owner)}\OperatorTok{;}

    \FunctionTok{payable}\NormalTok{(msg}\OperatorTok{.}\AttributeTok{sender}\NormalTok{)}\OperatorTok{.}\FunctionTok{transfer}\NormalTok{(}\FunctionTok{address}\NormalTok{(}\KeywordTok{this}\NormalTok{)}\OperatorTok{.}\AttributeTok{balance}\NormalTok{)}\OperatorTok{;}

\NormalTok{\}}
\end{Highlighting}
\end{Shaded}

\begin{itemize}
\tightlist
\item
  \textbf{Impacto}: Ataques mediante contratos intermediarios.
\item
  \textbf{Mitigación}: Usar msg.sender para autenticación
\end{itemize}

\subsubsection{ANEXO A.2.3. Inicialización insegura (contratos
upgradeables)}\label{anexo-a.2.3.-inicializaciuxf3n-insegura-contratos-upgradeables}

En contratos upgradeables, la inicialización se realiza mediante
funciones externas (initialize) en lugar de constructores. Si no están
protegidas, cualquier usuario puede ejecutarlas y asumir el control del
contrato.

Código vulnerable

\begin{Shaded}
\begin{Highlighting}[]
\KeywordTok{function} \FunctionTok{initialize}\NormalTok{(address \_owner) }\KeywordTok{public}\NormalTok{ \{}

\NormalTok{    owner }\OperatorTok{=}\NormalTok{ \_owner}\OperatorTok{;}

\NormalTok{\}}
\end{Highlighting}
\end{Shaded}

\begin{itemize}
\tightlist
\item
  \textbf{Impacto}: Un atacante puede inicializar el contrato antes que
  el legítimo propietario.
\item
  Mitigación:

  \begin{itemize}
  \tightlist
  \item
    Uso de
    \href{https://docs.openzeppelin.com/upgrades-plugins/writing-upgradeable}{initializer
    (OpenZeppelin)}
  \item
    Bloqueo de inicialización tras ejecución
  \end{itemize}
\end{itemize}

\subsection{ANEXO A.3. Vulnerabilidades económicas y del
entorno}\label{anexo-a.3.-vulnerabilidades-econuxf3micas-y-del-entorno}

\subsubsection{ANEXO A.3.1. Front-running /
MEV}\label{anexo-a.3.1.-front-running-mev}

Un atacante observa la mempool y ejecuta transacciones antes que la
víctimANEXO A.

Código vulnerable

\begin{Shaded}
\begin{Highlighting}[]
\KeywordTok{function} \FunctionTok{buy}\NormalTok{(uint price) }\KeywordTok{public}\NormalTok{ \{}

    \PreprocessorTok{require}\NormalTok{(price }\OperatorTok{==}\NormalTok{ currentPrice)}\OperatorTok{;}

    \CommentTok{// compra}

\NormalTok{\}}
\end{Highlighting}
\end{Shaded}

\begin{itemize}
\tightlist
\item
  \textbf{Impacto}: Manipulación de operaciones (arbitraje,
  liquidaciones, subastas).
\item
  Mitigación:

  \begin{itemize}
  \tightlist
  \item
    \href{https://medium.com/coinmonks/commit-reveal-scheme-in-solidity-c06eba4091bb}{\emph{Commit-reveal}}
  \item
    Subastas ciegas
  \item
    Uso de \emph{relayers} privados
  \end{itemize}
\end{itemize}

\subsubsection{ANEXO A.3.2. Dependencia de
oráculos}\label{anexo-a.3.2.-dependencia-de-oruxe1culos}

Uso de datos externos manipulables.

Código vulnerable

\begin{Shaded}
\begin{Highlighting}[]
\KeywordTok{function} \FunctionTok{getPrice}\NormalTok{() }\KeywordTok{public}\NormalTok{ view }\FunctionTok{returns}\NormalTok{ (uint) \{}

    \ControlFlowTok{return}\NormalTok{ externalOracle}\OperatorTok{.}\FunctionTok{price}\NormalTok{()}\OperatorTok{;}

\NormalTok{\}}
\end{Highlighting}
\end{Shaded}

\begin{itemize}
\tightlist
\item
  \textbf{Impacto}: Manipulación de precios en DeFi.
\item
  Mitigación:

  \begin{itemize}
  \tightlist
  \item
    Oráculos descentralizados (ej.
    \href{https://chain.link/}{Chainlink})
  \item
    Promedios temporales
    (\href{https://www.binance.com/es-MX/support/faq/detail/80655cc54d8a4b2bb8ea097001844fd1}{TWAP})
  \end{itemize}
\end{itemize}

\subsubsection{ANEXO A.3.3. Uso de
block.timestamp}\label{anexo-a.3.3.-uso-de-block.timestamp}

El uso de block.timestamp como fuente de aleatoriedad o para decisiones
críticas es inseguro, ya que su valor puede ser parcialmente manipulado
por mineros o validadores dentro de ciertos límites.

Código vulnerable

\begin{Shaded}
\begin{Highlighting}[]
\KeywordTok{function} \FunctionTok{random}\NormalTok{() }\KeywordTok{public}\NormalTok{ view }\FunctionTok{returns}\NormalTok{ (uint) \{}

    \ControlFlowTok{return} \FunctionTok{uint}\NormalTok{(}\FunctionTok{keccak256}\NormalTok{(abi}\OperatorTok{.}\FunctionTok{encodePacked}\NormalTok{(block}\OperatorTok{.}\AttributeTok{timestamp}\NormalTok{)))}\OperatorTok{;}

\NormalTok{\}}
\end{Highlighting}
\end{Shaded}

\begin{itemize}
\tightlist
\item
  \textbf{Impacto}: Resultados predecibles o manipulables por
  mineros/validadores.
\item
  Mitigación:

  \begin{itemize}
  \tightlist
  \item
    VRF
    (\href{https://chain.link/education-hub/verifiable-random-function-vrf}{Verifiable
    Random Functions})
  \item
    Fuentes externas verificables
  \end{itemize}
\end{itemize}

\subsection{ANEXO A.4. Errores lógicos de
negocio}\label{anexo-a.4.-errores-luxf3gicos-de-negocio}

\subsubsection{ANEXO A.4.1. Error en cálculo de
balances}\label{anexo-a.4.1.-error-en-cuxe1lculo-de-balances}

Errores en operadores lógicos o condiciones de validación pueden
provocar inconsistencias en la gestión de balances, especialmente en
casos límite donde las condiciones no cubren todos los escenarios
posibles.

Código vulnerable

\begin{Shaded}
\begin{Highlighting}[]
\KeywordTok{function} \FunctionTok{withdraw}\NormalTok{(uint amount) }\KeywordTok{public}\NormalTok{ \{}

    \PreprocessorTok{require}\NormalTok{(balances[msg}\OperatorTok{.}\AttributeTok{sender}\NormalTok{] }\OperatorTok{\textgreater{}}\NormalTok{ amount)}\OperatorTok{;}

\NormalTok{    balances[msg}\OperatorTok{.}\AttributeTok{sender}\NormalTok{] }\OperatorTok{{-}=}\NormalTok{ amount}\OperatorTok{;}

\NormalTok{\}}
\end{Highlighting}
\end{Shaded}

\begin{itemize}
\tightlist
\item
  \textbf{Impacto}: Comportamiento incorrecto en condiciones límite.
\item
  Mitigación:

  \begin{itemize}
  \tightlist
  \item
    Uso de \textgreater{} en lugar de \textgreater=.
  \end{itemize}
\end{itemize}

\subsubsection{ANEXO A.4.2. Distribución incorrecta de
recompensas}\label{anexo-a.4.2.-distribuciuxf3n-incorrecta-de-recompensas}

La lógica de distribución puede introducir errores debido a divisiones
enteras o falta de gestión de restos, provocando pérdida de precisión y
fondos no asignados correctamente.

Código vulnerable

\begin{Shaded}
\begin{Highlighting}[]
\KeywordTok{function} \FunctionTok{distribute}\NormalTok{() }\KeywordTok{public}\NormalTok{ \{}

\NormalTok{    uint reward }\OperatorTok{=}\NormalTok{ total }\OperatorTok{/}\NormalTok{ users}\OperatorTok{.}\AttributeTok{length}\OperatorTok{;}

    \ControlFlowTok{for}\NormalTok{ (uint i }\OperatorTok{=} \DecValTok{0}\OperatorTok{;}\NormalTok{ i }\OperatorTok{\textless{}}\NormalTok{ users}\OperatorTok{.}\AttributeTok{length}\OperatorTok{;}\NormalTok{ i}\OperatorTok{++}\NormalTok{) \{}

\NormalTok{        balances[users[i]] }\OperatorTok{+=}\NormalTok{ reward}\OperatorTok{;}

\NormalTok{    \}}

\NormalTok{\}}
\end{Highlighting}
\end{Shaded}

\begin{itemize}
\tightlist
\item
  Impacto:

  \begin{itemize}
  \tightlist
  \item
    Pérdida de fondos debido a errores de redondeo (truncamiento en
    división entera)
  \item
    Acumulación de saldo no distribuido en el contrato
  \item
    Distribuciones injustas entre usuarios
  \item
    Posibles vectores de explotación si un atacante manipula el número
    de participantes
  \end{itemize}
\item
  Mitigación:

  \begin{itemize}
  \tightlist
  \item
    Uso de patrones de distribución que gestionen residuos (por ejemplo,
    acumuladores o ``\emph{remainder handling}'')
  \item
    Empleo de mayor precisión mediante escalado
    (\href{https://rareskills.io/post/solidity-fixed-point}{\emph{fixed-point
    arithmetic}})
  \item
    Validación de invariantes económicas (la suma distribuida debe
    coincidir con el total)
  \item
    Testing específico de casos límite (número de usuarios, valores
    pequeños, etc.)
  \end{itemize}
\end{itemize}

\subsubsection{ANEXO A.4.3. Estados
inconsistentes}\label{anexo-a.4.3.-estados-inconsistentes}

La falta de control adecuado sobre las transiciones de estado puede
permitir la ejecución de funciones en condiciones no válidas, generando
comportamientos inconsistentes en el contrato.

Código vulnerable

\begin{Shaded}
\begin{Highlighting}[]
\KeywordTok{enum}\NormalTok{ State \{ Open}\OperatorTok{,}\NormalTok{ Closed \}}

\NormalTok{State }\KeywordTok{public}\NormalTok{ state}\OperatorTok{;}

\KeywordTok{function} \FunctionTok{close}\NormalTok{() }\KeywordTok{public}\NormalTok{ \{}

\NormalTok{    state }\OperatorTok{=}\NormalTok{ State}\OperatorTok{.}\AttributeTok{Closed}\OperatorTok{;}

\NormalTok{\}}

\KeywordTok{function} \FunctionTok{bid}\NormalTok{() }\KeywordTok{public}\NormalTok{ payable \{}

    \PreprocessorTok{require}\NormalTok{(state }\OperatorTok{==}\NormalTok{ State}\OperatorTok{.}\AttributeTok{Open}\NormalTok{)}\OperatorTok{;}

\NormalTok{\}}
\end{Highlighting}
\end{Shaded}

\begin{itemize}
\tightlist
\item
  Impacto:

  \begin{itemize}
  \tightlist
  \item
    Ejecución de funciones en estados no válidos
  \item
    Comportamiento inesperado del contrato
  \item
    Bloqueo o bypass de lógica de negocio
  \item
    Posible explotación combinada con otras vulnerabilidades (por
    ejemplo, \emph{front-running} o \emph{reentrancy}
  \end{itemize}
\item
  Mitigación:

  \begin{itemize}
  \tightlist
  \item
    Implementación de máquinas de estados explícitas y completas
  \item
    Uso de modificadores para validar estado (inState(State.Open))
  \item
    Restricción de transiciones de estado válidas
  \item
    Aplicación de patrones
    \href{https://fravoll.github.io/solidity-patterns/state_machine.html}{\emph{state
    machines}}
  \end{itemize}
\end{itemize}

\subsection{ANEXO A.5. Resumen de
vulnerabilidades}\label{anexo-a.5.-resumen-de-vulnerabilidades}

{\def\LTcaptype{none} % do not increment counter
\begin{longtable}[]{@{}
  >{\raggedright\arraybackslash}p{(\linewidth - 12\tabcolsep) * \real{0.0661}}
  >{\raggedright\arraybackslash}p{(\linewidth - 12\tabcolsep) * \real{0.1074}}
  >{\raggedright\arraybackslash}p{(\linewidth - 12\tabcolsep) * \real{0.1736}}
  >{\raggedright\arraybackslash}p{(\linewidth - 12\tabcolsep) * \real{0.1570}}
  >{\raggedright\arraybackslash}p{(\linewidth - 12\tabcolsep) * \real{0.0909}}
  >{\raggedright\arraybackslash}p{(\linewidth - 12\tabcolsep) * \real{0.2479}}
  >{\raggedright\arraybackslash}p{(\linewidth - 12\tabcolsep) * \real{0.1570}}@{}}
\toprule\noalign{}
\begin{minipage}[b]{\linewidth}\raggedright
\textbf{\emph{ID}}
\end{minipage} & \begin{minipage}[b]{\linewidth}\raggedright
\textbf{Categoría}
\end{minipage} & \begin{minipage}[b]{\linewidth}\raggedright
\textbf{Vulnerabilidad}
\end{minipage} & \begin{minipage}[b]{\linewidth}\raggedright
\textbf{Tipo}
\end{minipage} & \begin{minipage}[b]{\linewidth}\raggedright
\textbf{Impacto}
\end{minipage} & \begin{minipage}[b]{\linewidth}\raggedright
\textbf{Detectable automáticamente}
\end{minipage} & \begin{minipage}[b]{\linewidth}\raggedright
\textbf{Ejemplo sección}
\end{minipage} \\
\midrule\noalign{}
\endhead
\bottomrule\noalign{}
\endlastfoot
\emph{V1} & Técnica & Reentrancy & Ejecución & Crítico & Sí
(Slither/Mythril) & ANEXO A.1.1 \\
\emph{V2} & Técnica & Overflow & Aritmético & Medio & Sí (Slither) &
ANEXO A.1.2 \\
\emph{V3} & Técnica & Delegatecall & Ejecución & Crítico & Parcial &
ANEXO A.1.3 \\
\emph{V4} & Control & Acceso no restringido & Autorización & Crítico &
Sí & ANEXO A.2.1 \\
\emph{V5} & Control & tx.origin & Autenticación & Alto & Sí & ANEXO
A.2.2 \\
\emph{V6} & Económica & Front-running & MEV & Alto & No & ANEXO A.3.1 \\
\emph{V7} & Económica & Oráculos & Dependencia externa & Crítico & No &
ANEXO A.3.2 \\
\emph{V8} & Lógica & Error de balance & Lógica & Variable & No & ANEXO
A.4.1 \\
\end{longtable}
}


\chapter{Entorno virtual Python}
\section{ANEXO B. Entorno virtual
Python}\label{anexo-b.-entorno-virtual-python}

El desarrollo de una librería Python destinada a integrar múltiples
herramientas de análisis de seguridad plantea, desde el inicio, un
problema de gestión de dependencias que no debe subestimarse. Las
herramientas que se integran en la solución propuesta, como Slither,
Mythril o Echidna, tienen requisitos de versión específicos y en
ocasiones incompatibles entre sí cuando se instalan en el entorno global
del sistema. Esta situación, conocida en el ecosistema Python como
dependency hell, puede provocar conflictos silenciosos difíciles de
depurar y comprometer la reproducibilidad del entorno de desarrollo.

Un entorno virtual (virtual environment) es un directorio aislado que
contiene una instalación independiente del intérprete Python junto con
sus propios paquetes y dependencias, sin interferir con el sistema
global ni con otros proyectos. Esta separación proporciona varias
ventajas fundamentales en el contexto del presente trabajo:

En primer lugar, garantiza el \textbf{aislamiento de dependencias}, de
forma que cada proyecto mantiene sus propias versiones de bibliotecas
sin afectar al resto del sistema. Esto resulta especialmente relevante
cuando diferentes herramientas de análisis requieren versiones distintas
de una misma dependencia transitiva.

En segundo lugar, favorece la \textbf{reproducibilidad del entorno de
desarrollo}. Todos los integrantes del equipo pueden trabajar
exactamente con las mismas versiones de todas las dependencias,
eliminando la variabilidad asociada a instalaciones manuales y
garantizando que los resultados obtenidos durante el desarrollo son
consistentes independientemente del sistema operativo o configuración
personal de cada desarrollador.

En tercer lugar, facilita el \textbf{ciclo de vida} del proyecto al
delimitar claramente qué paquetes pertenecen al proyecto y cuáles son
del sistema, simplificando tanto la distribución de la librería como su
posterior publicación en registros públicos como PyPI.

\subsection{ANEXO B.1. Tipos de entornos
virtuales}\label{anexo-b.1.-tipos-de-entornos-virtuales}

En el ecosistema Python existen varias alternativas para la gestión de
entornos virtuales, con distintos niveles de abstracción y
funcionalidad.

El módulo estándar \texttt{venv}, incluido en la biblioteca estándar
desde Python 3.3, permite crear entornos virtuales básicos mediante el
comando \texttt{python\ -m\ venv\ .venv}. Sin embargo, este enfoque no
incluye gestión de dependencias ni ficheros de bloqueo
(\emph{lockfiles}), por lo que debe complementarse con herramientas
adicionales como \texttt{pip} y \texttt{pip-tools}.

\texttt{virtualenv} es una alternativa anterior al módulo estándar, con
mayor compatibilidad con versiones antiguas de Python y algunas
funcionalidades adicionales, aunque en la práctica ha quedado desplazada
por las herramientas modernas.

\texttt{conda} ofrece un modelo más completo que combina gestión de
entornos con gestión de paquetes, incluyendo dependencias no Python. Es
habitual en entornos científicos y de análisis de datos, pero introduce
una complejidad y un tamaño innecesarios para un proyecto centrado en el
ecosistema Python puro.

Herramientas como \texttt{poetry} o \texttt{pipenv} representan un nivel
superior de abstracción, combinando la gestión de entornos virtuales con
la resolución de dependencias, la generación de ficheros de bloqueo y el
ciclo de publicación de paquetes. Su adopción en proyectos profesionales
se ha generalizado en los últimos años.

Finalmente, \texttt{uv} constituye la herramienta de última generación
en este espacio, combinando todas las funcionalidades anteriores en una
solución de rendimiento muy superior, como se detalla en la sección
siguiente.

\subsection{ANEXO B.2. Herramienta a usar:
uv}\label{anexo-b.2.-herramienta-a-usar-uv}

\texttt{uv} es un gestor de paquetes y proyectos Python de alto
rendimiento desarrollado por Astral, la empresa creadora del formateador
Ruff. Implementado en Rust, se presenta como una herramienta unificada
capaz de sustituir a \texttt{pip}, \texttt{pip-tools}, \texttt{pipx},
\texttt{poetry}, \texttt{pyenv}, \texttt{twine} y \texttt{virtualenv}
mediante una interfaz de línea de comandos coherente. Su desarrollo se
orienta tanto a la velocidad de ejecución como a la corrección en la
resolución de dependencias y a la facilidad de adopción en proyectos de
diferente escala y complejidad.

\subsubsection{ANEXO B.2.1. Ventajas de
uv}\label{anexo-b.2.1.-ventajas-de-uv}

La principal ventaja diferencial de \texttt{uv} respecto a las
alternativas existentes es su rendimiento. Según los propios
\emph{benchmarks} publicados en su documentación oficial, \texttt{uv} es
entre 10 y 100 veces más rápido que \texttt{pip} en operaciones de
instalación de paquetes con caché caliente. Esta diferencia resulta
perceptible en la práctica, especialmente durante las fases de
incorporación de nuevos integrantes al equipo o en entornos de
integración continua donde el entorno debe reconstruirse frecuentemente.

Más allá del rendimiento, \texttt{uv} ofrece un conjunto de ventajas
relevantes para el desarrollo de la librería propuesta en este trabajo.
Gestiona automáticamente los entornos virtuales asociados a cada
proyecto, sin necesidad de crearlos ni activarlos manualmente. Genera y
mantiene un fichero de bloqueo universal (\texttt{uv.lock}) que
garantiza instalaciones reproducibles en cualquier plataforma. Permite
gestionar múltiples versiones del intérprete Python e instalar la
versión adecuada de forma automática si no está disponible en el
sistema. Además, su diseño es compatible con los estándares del
ecosistema Python (\texttt{pyproject.toml}, PEP 517, PEP 621), lo que
facilita la integración con otras herramientas y la publicación en
registros de paquetes.

\subsubsection{ANEXO B.2.2. Qué proporciona uv en el contexto de este
proyecto}\label{anexo-b.2.2.-quuxe9-proporciona-uv-en-el-contexto-de-este-proyecto}

Para el desarrollo de la librería propuesta, uv proporciona un conjunto
de funcionalidades que cubren todo el ciclo de vida del proyecto, desde
la inicialización hasta la publicación.

\textbf{Gestión de dependencias y sincronización del entorno.} Una vez
definidas las dependencias del proyecto en el fichero
\texttt{pyproject.toml}, cualquier integrante del equipo puede
reproducir exactamente el mismo entorno ejecutando un único comando:

\begin{Shaded}
\begin{Highlighting}[]
\ExtensionTok{uv}\NormalTok{ sync}
\end{Highlighting}
\end{Shaded}

Este comando resuelve las dependencias declaradas, instala las versiones
exactas registradas en el fichero de bloqueo \texttt{uv.lock} y
configura el entorno virtual del proyecto de forma automática. La
simplicidad de este flujo elimina los problemas habituales de
divergencia entre entornos de desarrollo individuales, garantizando que
todos los miembros del equipo trabajan con las mismas versiones de
Slither, Mythril, Echidna y el resto de dependencias de la librería.

\textbf{Inicialización de proyectos.} \texttt{uv} proporciona soporte
integrado para la creación de nuevos proyectos mediante el comando
\texttt{uv\ init}. Para el caso específico de una librería Python
destinada a ser importada por otros proyectos o publicada en PyPI, se
utiliza la opción \texttt{-\/-lib}:

\begin{Shaded}
\begin{Highlighting}[]
\ExtensionTok{uv}\NormalTok{ init }\AttributeTok{{-}{-}lib}\NormalTok{ evmaudit}
\end{Highlighting}
\end{Shaded}

Este comando genera automáticamente la estructura de directorios
recomendada para una librería Python, incluyendo el fichero
\texttt{pyproject.toml} con los metadatos del proyecto, el directorio
\texttt{src/} con el paquete principal y los ficheros de configuración
necesarios para la construcción y distribución. El uso de la disposición
\texttt{src/} (\emph{src layout}) es la práctica recomendada actualmente
para proyectos publicables, ya que evita problemas habituales
relacionados con la importación del paquete desde el directorio raíz
durante el desarrollo.

\textbf{Construcción de distribuciones.} \texttt{uv} integra soporte
nativo para la generación de distribuciones instalables mediante el
comando \texttt{uv\ build}, que produce tanto el archivo fuente
(\emph{sdist}) como la rueda binaria (\emph{wheel}) del paquete:

\begin{Shaded}
\begin{Highlighting}[]
\ExtensionTok{uv}\NormalTok{ build}
\end{Highlighting}
\end{Shaded}

El resultado son los artefactos estándar de distribución Python ubicados
en el directorio \texttt{dist/}, listos para ser publicados o
distribuidos directamente.

\textbf{Publicación de paquetes.} El ciclo se completa con el soporte
para publicación en registros de paquetes, incluyendo PyPI, mediante el
comando \texttt{uv\ publish}:

\begin{Shaded}
\begin{Highlighting}[]
\ExtensionTok{uv}\NormalTok{ publish}
\end{Highlighting}
\end{Shaded}

Este comando gestiona la autenticación y la subida de los artefactos
generados, cubriendo el flujo completo que anteriormente requería
herramientas adicionales como \texttt{twine}.

En conjunto, \texttt{uv} unifica en una sola herramienta todo el ciclo
de vida del proyecto: inicialización, gestión de dependencias,
sincronización del entorno, construcción de distribuciones y
publicación. Esta integración reduce la fricción en el desarrollo
colaborativo y facilita la adopción de prácticas profesionales de
gestión de proyectos Python desde las primeras fases del trabajo.


\chapter{Distribución mediante PyPI}
\section{ANEXO C. DISTRIBUCIÓN Y PUBLICACIÓN EN EL REGISTRO DE PAQUETES
PYPI}\label{anexo-c.-distribuciuxf3n-y-publicaciuxf3n-en-el-registro-de-paquetes-pypi}

El ciclo de desarrollo de la librería propuesta culmina con su fase de
distribución, permitiendo que las herramientas de análisis de seguridad
implementadas sean accesibles e integrables por la comunidad de
desarrollo y auditoría de \emph{smart contracts}. Para asegurar una
distribución estandarizada y eficiente dentro del ecosistema Python, se
ha seleccionado el índice oficial de paquetes PyPI (\emph{Python Package
Index}). La gestión de este proceso se unifica bajo la herramienta
\texttt{uv}, garantizando la consistencia desde la compilación de los
artefactos hasta su publicación definitiva. \#\# ANEXO C.1.
Configuración de Metadatos del Proyecto (pyproject.toml)

El paso previo indispensable para la distribución consiste en la
definición inequívoca de los metadatos y la especificación del sistema
de construcción (\emph{build system}) en el archivo de configuración
\texttt{pyproject.toml}, ubicado en la raíz del paquete. Este
procedimiento se rige bajo los estándares modernos de empaquetado de
Python (PEP 517 y PEP 621).

Para este proyecto, se ha adoptado \texttt{hatchling} como \emph{build
backend}, debido a su ligereza, velocidad y compatibilidad nativa con
las especificaciones del ecosistema actual. A continuación, se presenta
la estructura de configuración requerida para delimitar las propiedades
de la librería \texttt{evmaudit}:

\begin{Shaded}
\begin{Highlighting}[]
\KeywordTok{[project]}
\DataTypeTok{name} \OperatorTok{=} \StringTok{"evmaudit"}
\DataTypeTok{version} \OperatorTok{=} \StringTok{"0.1.0"}
\DataTypeTok{description} \OperatorTok{=} \StringTok{"Add your description here"}
\DataTypeTok{readme} \OperatorTok{=} \StringTok{"README.md"}
\DataTypeTok{authors} \OperatorTok{=} \OperatorTok{[}
    \OperatorTok{\{ }\DataTypeTok{name}\OperatorTok{ =} \StringTok{"Daniel Rovira Martínez"}\OperatorTok{, }\DataTypeTok{email}\OperatorTok{ =} \StringTok{"pardalaco@gmail.com"}\OperatorTok{ \},}
    \OperatorTok{\{ }\DataTypeTok{name}\OperatorTok{ =} \StringTok{"Paula Suárez Prieto"}\OperatorTok{, }\DataTypeTok{email}\OperatorTok{ =} \StringTok{"Paula Suárez Prieto"}\OperatorTok{ \},}
    \OperatorTok{\{ }\DataTypeTok{name}\OperatorTok{ =} \StringTok{"Adrián Moreno Martín"}\OperatorTok{, }\DataTypeTok{email}\OperatorTok{ =} \StringTok{"adrimore2@gmail.com"}\OperatorTok{ \}}
\OperatorTok{]}
\DataTypeTok{requires{-}python} \OperatorTok{=} \StringTok{"\textgreater{}=3.12"}
\DataTypeTok{dependencies} \OperatorTok{=} \OperatorTok{[}
    \StringTok{"eth\textgreater{}=0.0.1"}\OperatorTok{,}
    \StringTok{"mythril\textgreater{}=0.24.8"}\OperatorTok{,}
    \StringTok{"setuptools\textless{}70.0.0"}\OperatorTok{,}
    \StringTok{"slither{-}analyzer\textgreater{}=0.9.2"}\OperatorTok{,}
    \StringTok{"solc{-}select\textgreater{}=1.2.0"}\OperatorTok{,]}

\KeywordTok{[project.scripts]}
\DataTypeTok{evmaudit} \OperatorTok{=} \StringTok{"evmaudit.main:main"}

\KeywordTok{[build{-}system]}
\DataTypeTok{requires} \OperatorTok{=} \OperatorTok{[}\StringTok{"uv\_build\textgreater{}=0.11.15,\textless{}0.12.0"}\OperatorTok{]}
\DataTypeTok{build{-}backend} \OperatorTok{=} \StringTok{"uv\_build"}
\end{Highlighting}
\end{Shaded}

Los clasificadores (\emph{classifiers}) incluidos permiten categorizar
la librería dentro del índice público, facilitando su indexación en
función de la licencia de código abierto seleccionada, la compatibilidad
del sistema operativo y las versiones soportadas del intérprete Python.

\subsection{ANEXO C.2. Proceso de Compilación del
Paquete}\label{anexo-c.2.-proceso-de-compilaciuxf3n-del-paquete}

Una vez validados los metadatos, se procede a la generación de los
archivos de distribución. Este proceso transforma el código fuente
estructurado en el directorio \texttt{src/} en artefactos instalables e
independientes del entorno de desarrollo.

La ejecución del comando unificado de compilación abstrae la complejidad
de las herramientas tradicionales:

\begin{Shaded}
\begin{Highlighting}[]
\ExtensionTok{uv}\NormalTok{ build}
\end{Highlighting}
\end{Shaded}

Este comando genera de forma nativa dos tipos de distribuciones estándar
dentro del directorio \texttt{dist/}:

\begin{itemize}
\tightlist
\item
  \textbf{Distribución de código fuente (\emph{sdist} o}
  \texttt{.tar.gz}\textbf{):} Un archivo comprimido que contiene el
  código fuente original y los archivos de configuración, actuando como
  respaldo para plataformas o configuraciones no previstas.
\item
  \textbf{Distribución binaria compilada (\emph{wheel} o}
  \texttt{.whl}\textbf{):} El formato de empaquetado moderno que permite
  una instalación directa y optimizada en el sistema destino, evitando
  la necesidad de compilar el código en la máquina del usuario final.
\end{itemize}

En contextos de desarrollo complejos donde el paquete forma parte de un
repositorio principal o arquitectura de \emph{workspace} (como el
entorno de trabajo \texttt{TFM-UNIR}), \texttt{uv} permite gestionar la
compilación de manera localizada. Para forzar la construcción exclusiva
del subproyecto desde su propio directorio y evitar colisiones en la
raíz global, se aplica la opción de empaquetado específico:

\begin{Shaded}
\begin{Highlighting}[]
    \ExtensionTok{uv}\NormalTok{ build }\AttributeTok{{-}{-}package}\NormalTok{ evmaudit}
\end{Highlighting}
\end{Shaded}

\subsection{ANEXO C.3. Seguridad y Autenticación en la
Publicación}\label{anexo-c.3.-seguridad-y-autenticaciuxf3n-en-la-publicaciuxf3n}

La publicación de código en repositorios públicos exige mecanismos
estrictos de control de acceso para prevenir vectores de ataque basados
en la cadena de suministro (\emph{supply chain attacks}). Por razones de
seguridad, se desestima el uso de contraseñas de usuario tradicionales
en favor de la autenticación basada en \textbf{Tokens de API.}

El proceso de despliegue requiere la obtención de un \emph{token} con
prefijo \texttt{pypi-} generado desde el panel de control de PyPI. En la
primera interacción, el alcance del \emph{token} se configura de manera
global; no obstante, una vez realizada la primera subida con éxito, la
buena práctica metodológica dicta restringir los permisos del token de
manera exclusiva al ámbito del paquete \texttt{evmaudit}, minimizando
así la superficie de exposición en caso de compromiso de la credencial.

\subsection{ANEXO C.4. Ejecución del
Despliegue}\label{anexo-c.4.-ejecuciuxf3n-del-despliegue}

Con los artefactos ubicados en el directorio \texttt{dist/} y las
credenciales expedidas, se procede a la transferencia segura hacia los
servidores de PyPI. El comando \texttt{uv\ publish} automatiza la
verificación de integridad mediante \emph{hashes} criptográficos y
realiza la subida en un único paso:

\begin{Shaded}
\begin{Highlighting}[]
\ExtensionTok{uv}\NormalTok{ publish}
\end{Highlighting}
\end{Shaded}

Durante el flujo interactivo en la línea de comandos, el sistema
requiere la introducción del identificador genérico
\texttt{\_\_token\_\_} en el campo de usuario, seguido de la clave
alfanumérica del token de API en el campo de contraseña. Con el objetivo
de optimizar este flujo en entornos de Integración Continua (CI/CD) o
evitar la inserción manual recurrente, es posible exportar temporalmente
la credencial en el entorno de la terminal actual:

\begin{Shaded}
\begin{Highlighting}[]
\BuiltInTok{export} \VariableTok{UV\_PUBLISH\_TOKEN}\OperatorTok{=}\StringTok{"pypi{-}tu{-}token{-}aqui"}
\ExtensionTok{uv}\NormalTok{ publish}
\end{Highlighting}
\end{Shaded}

Tras la finalización exitosa del proceso, el paquete queda registrado
globalmente, permitiendo su incorporación inmediata en otros proyectos
mediante los gestores tradicionales del ecosistema:

\begin{Shaded}
\begin{Highlighting}[]
\ExtensionTok{pip}\NormalTok{ install evmaudit}
\end{Highlighting}
\end{Shaded}

O bien, aprovechando los beneficios de rendimiento de la herramienta
unificada del proyecto:

\begin{Shaded}
\begin{Highlighting}[]
\ExtensionTok{uv}\NormalTok{ add evmaudit}
\end{Highlighting}
\end{Shaded}

\subsection{ANEXO C.5. Ciclo de Mantenimiento y Actualización de
Versiones}\label{anexo-c.5.-ciclo-de-mantenimiento-y-actualizaciuxf3n-de-versiones}

La evolución de la librería para la corrección de vulnerabilidades o la
integración de nuevas capacidades de análisis requiere una política
estricta de control de versiones. El flujo metodológico establecido para
la liberación de actualizaciones iterativas consta de tres fases
secuenciales: - \textbf{Incremento del número de versión:} Modificación
manual del campo \texttt{version} en el archivo \texttt{pyproject.toml}
siguiendo el estándar de Versionado Semántico (ej. de \texttt{0.1.0} a
\texttt{0.1.1}). - \textbf{Saneamiento del directorio de distribución:}
Eliminación de los artefactos obsoletos del directorio \texttt{dist/}
para mitigar el riesgo de duplicidad o subidas erróneas de versiones
previas. - \textbf{Reconstrucción y despliegue:} Ejecución consecutiva
de los procesos de empaquetado y transferencia:

\begin{Shaded}
\begin{Highlighting}[]
\ExtensionTok{uv}\NormalTok{ build}
\ExtensionTok{uv}\NormalTok{ publish}
\end{Highlighting}
\end{Shaded}

Esta sistemática asegura que cada iteración de la herramienta de
auditoría de la EVM mantenga la trazabilidad, la coherencia histórica y
la disponibilidad pública necesarias para un entorno de producción
académica y profesional.


\chapter{Docker}
\section{Anexo D Docker}\label{anexo-d-docker}

\subsection{ANEXO D.1. CONTENEDORIZACIÓN E INFRAESTRUCTURA DE DESPLIEGUE
(DOCKER)}\label{anexo-d.1.-contenedorizaciuxf3n-e-infraestructura-de-despliegue-docker}

En el ámbito del desarrollo de software moderno y la ciberseguridad, la
reproducibilidad del entorno de ejecución constituye un pilar crítico.
Tradicionalmente, el despliegue de aplicaciones que integran múltiples
herramientas de análisis (como compiladores de Solidity y motores de
ejecución simbólica) se enfrentaba al problema de ``funciona en mi
máquina'', derivado de las discrepancias en las versiones de las
dependencias, librerías del sistema operativo y configuraciones locales.
Para mitigar este riesgo, el presente proyecto adopta una arquitectura
basada en contenedores de aplicación a través del ecosistema de
\textbf{Docker} y \textbf{Docker Compose}.

\subsection{ANEXO D.2. Fundamentos de Contenedorización: Docker y Docker
Compose}\label{anexo-d.2.-fundamentos-de-contenedorizaciuxf3n-docker-y-docker-compose}

Docker es una plataforma de código abierto basada en la tecnología de
contenedorización, la cual permite empaquetar una aplicación y todas sus
dependencias (binarios, librerías, archivos de configuración) en una
unidad estandarizada denominada \textbf{contenedor}. A diferencia de la
virtualización tradicional, Docker opera mediante la virtualización a
nivel de sistema operativo, compartiendo el núcleo (kernel) del sistema
anfitrión pero ejecutando los procesos en espacios de usuario
completamente aislados a través de namespaces y cgroups. Desde la
perspectiva de la seguridad, este aislamiento garantiza que los procesos
del pipeline de auditoría de EVMAudit se ejecuten de forma confinada,
mitigando el impacto en la infraestructura anfitriona ante la eventual
ejecución de código arbitrario o inesperado durante el análisis de
contratos inteligentes.

Por su parte, \textbf{Docker Compose} es la herramienta diseñada para
definir y orquestar aplicaciones Docker multi-contenedor. Mediante el
uso de un archivo de configuración declarativo en formato YAML
(docker-compose.yml), permite definir con precisión los servicios que
componen el sistema, sus dependencias de arranque, la exposición de
puertos hacia el exterior, la creación de redes aisladas y la asignación
de volúmenes persistentes. En el contexto de EVMAudit, actúa como el
motor de despliegue unificado, permitiendo al administrador inicializar
toda la infraestructura del TFM (servidor FastAPI, interfaz web y
almacenamiento de informes) de manera centralizada.

\subsection{ANEXO D.3. Estrategia de Construcción de la Imagen
(Dockerfile)}\label{anexo-d.3.-estrategia-de-construcciuxf3n-de-la-imagen-dockerfile}

La construcción de la imagen se define en un único Dockerfile
optimizado. Debido a que el pipeline de análisis requiere interactuar
con el sistema operativo para invocar compiladores y binarios de
seguridad, se ha seleccionado \textbf{Ubuntu 22.04} como imagen base,
proporcionando un entorno estable y con soporte extendido para
dependencias nativas de Linux en arquitectura amd64.

El proceso de aprovisionamiento de la imagen se divide en las siguientes
fases críticas: 1. \textbf{Entorno y Variables de Sistema}: Se
configuran las variables de entorno \texttt{PYTHONDONTWRITEBYTECODE=1} y
\texttt{PYTHONUNBUFFERED=1} para optimizar la ejecución de Python dentro
del contenedor, evitando la escritura de residuos binarios y forzando el
volcado de logs en tiempo real. Asimismo, se establece
\texttt{DEBIAN\_FRONTEND=noninteractive} para suprimir diálogos
interactivos durante la instalación de paquetes. 2.
\textbf{Aprovisionamiento de Compiladores (Solidity):} Se añade el
repositorio PPA oficial de Ethereum (ppa:ethereum/ethereum) para
incorporar el compilador nativo de Solidity (solc). Posteriormente, se
instala la utilidad \texttt{solc-select} mediante el gestor de paquetes
de Python para automatizar la descarga y conmutación de versiones. 3.
\textbf{Integración del Fuzzer Echidna:} Dado que Echidna se distribuye
de manera óptima como un binario estático para Linux, el contenedor
automatiza su descarga directa (versión v2.3.2) desde los repositorios
oficiales de \emph{Crytic}, procediendo a su extracción e instalación en
\texttt{/usr/local/bin/} para garantizar su disponibilidad inmediata en
el PATH del sistema. 4. \textbf{Optimización de Dependencias con
\texttt{uv} (Multi-stage Build):} Con el objetivo de minimizar los
tiempos de construcción y asegurar una gestión eficiente de los paquetes
de Python, se emplea un mecanismo de construcción en etapas múltiples
(Multi-stage build), importando los binarios optimizados del gestor uv
directamente desde su imagen oficial en el registro de GitHub
(ghcr.io/astral-sh/uv:latest). 5. \textbf{Instalación del Paquete
Local:} Tras establecer el directorio de trabajo en /app y copiar el
código fuente , se ejecuta el comando
\texttt{uv\ sync\ -\/-frozen\ -\/-no-cache}. Esto resuelve de forma
determinista el grafo de dependencias del archivo uv.lock, registrando
el paquete local editable evmaudit e instalando el servidor ASGI Uvicorn
sin almacenar datos residuales en la caché de la imagen.

\subsection{ANEXO D.4. Orquestación de Servicios (Docker
Compose)}\label{anexo-d.4.-orquestaciuxf3n-de-servicios-docker-compose}

La coordinación del contenedor web y sus dependencias con el sistema
anfitrión se gestiona de forma declarativa mediante un archivo
docker-compose.yml. La especificación del servicio, denominado
evmaudit-web, se fundamenta en tres pilares de ingeniería:

\begin{itemize}
\tightlist
\item
  \textbf{Persistencia de Datos mediante Volúmenes}: Con el fin de dotar
  al sistema de un estado persistente (pese a la naturaleza efímera de
  los contenedores), se realiza un mapeo directo de directorios del host
  hacia el contenedor:
\item
  \textbf{\texttt{./jsons/\_uploads:/app/jsons/\_uploads}}: Almacena de
  forma persistente los contratos Solidity cargados por los usuarios,
  los wrappers intermedios generados para Echidna y los informes de
  auditoría finales en formato JSON y Markdown.
\item
  \textbf{\texttt{./contracts:/app/contracts}}: Habilita un volumen
  opcional para la auditoría directa de Smart Contracts locales sin
  necesidad de interactuar con la interfaz web.
\item
  \textbf{Aislamiento de Red y Mapeo de Puertos}: Se expone el puerto
  8080 del contenedor hacia el puerto 8080 del sistema anfitrión. Esto
  permite redirigir el tráfico HTTP de la interfaz construida en
  HTML5/Vanilla JS hacia el backend desarrollado en FastAPI de forma
  transparente.
\item
  \textbf{Tolerancia a Fallos}: Se implementa la política de reinicio
  restart: unless-stopped. Esta directiva asegura la alta disponibilidad
  del servicio ante excepciones imprevistas en el motor de ejecución
  simbólica (Mythril) o caídas del propio demonio de Docker,
  garantizando que el servicio web vuelva a levantarse de manera
  automática salvo interrupción explícita del administrador.
\end{itemize}

\subsection{ANEXO D.5. Flujo de Despliegue y Ciclo de Vida del
Contenedor}\label{anexo-d.5.-flujo-de-despliegue-y-ciclo-de-vida-del-contenedor}

Para la puesta en marcha de la infraestructura local en entornos de
desarrollo o evaluación, el ciclo de vida del contenedor se administra
mediante el estándar de comandos de Docker Compose: 1. \textbf{Fase de
Construcción (Build):} Compilación de la imagen e instalación del
entorno virtual determinista:

\begin{Shaded}
\begin{Highlighting}[]
\ExtensionTok{docker}\NormalTok{ compose build}
\end{Highlighting}
\end{Shaded}

\begin{enumerate}
\def\labelenumi{\arabic{enumi}.}
\setcounter{enumi}{1}
\tightlist
\item
  \textbf{Fase de Inicialización (Up):} Despliegue e instanciación del
  servicio en segundo plano (detached mode):
\end{enumerate}

\begin{Shaded}
\begin{Highlighting}[]
\ExtensionTok{docker}\NormalTok{ compose up }\AttributeTok{{-}d}
\end{Highlighting}
\end{Shaded}

\begin{enumerate}
\def\labelenumi{\arabic{enumi}.}
\setcounter{enumi}{2}
\tightlist
\item
  \textbf{Fase de Auditoría de Ejecución (Logs):} Inspección de la
  salida estándar del contenedor para la monitorización de los análisis
  en curso:
\end{enumerate}

\begin{Shaded}
\begin{Highlighting}[]
\ExtensionTok{docker}\NormalTok{ compose logs }\AttributeTok{{-}f}\NormalTok{ evmaudit{-}web}
\end{Highlighting}
\end{Shaded}

\begin{enumerate}
\def\labelenumi{\arabic{enumi}.}
\setcounter{enumi}{3}
\tightlist
\item
  \textbf{Fase de Parada (Down):} Interrupción y eliminación de los
  contenedores activos salvaguardando la integridad de los datos de las
  auditorías gracias a los volúmenes enlazados: docker compose down
\end{enumerate}


\chapter{Aplicación web}
\input{capitulos/Z3 - Anexo E Aplicación web}

\chapter{CI/CD}
\section{ANEXO F. ENTORNO DE INTEGRACIÓN CONTINUA (CI/CD) Y PUBLICACIÓN
AUTOMATIZADA}\label{anexo-f.-entorno-de-integraciuxf3n-continua-cicd-y-publicaciuxf3n-automatizada}

Para garantizar la integridad del software durante el ciclo de vida del
desarrollo y agilizar el flujo de despliegue, se ha diseñado e
implementado un pipeline de Integración Continua (CI) basado en
\textbf{GitHub Actions}. Esta estrategia de ingeniería de software
permite automatizar la compilación, verificación y empaquetado de la
aplicación en cada iteración, mitigando los riesgos asociados a la
integración manual de código y asegurando la disponibilidad inmediata de
artefactos listos para producción.

\subsection{ANEXO F.1. Arquitectura del Workflow y Disparadores
(Triggers)}\label{anexo-f.1.-arquitectura-del-workflow-y-disparadores-triggers}

El flujo de trabajo automatizado se define de manera declarativa
mediante la sintaxis YAML de GitHub Actions. Con el propósito de
optimizar los recursos de cómputo y mantener un control estricto sobre
la estabilidad de la rama principal, se han configurado dos disparadores
específicos:

\begin{itemize}
\item
  \textbf{Eventos de Empuje (\emph{Push}):} El pipeline se ejecuta de
  forma automática ante cualquier consolidación directa de código en la
  rama \texttt{main}, asegurando que cada incremento de software sea
  evaluado y empaquetado de inmediato.
\item
  \textbf{Solicitudes de Extracción (\emph{Pull Requests}):} La
  automatización actúa como una barrera de calidad (\emph{quality gate})
  ante cualquier intento de fusión hacia la rama \texttt{main}. Esto
  permite validar que las modificaciones propuestas por los
  desarrolladores no rompan el proceso de construcción del contenedor
  antes de que el código sea integrado definitivamente.
\end{itemize}

El entorno de ejecución seleccionado para los trabajos (\emph{jobs}) es
\texttt{ubuntu-latest}, lo que proporciona un entorno virtual limpio,
aislado y actualizado de manera nativa por la infraestructura de GitHub.

\subsection{ANEXO F.2. Desglose Técnico de las Etapas del
Pipeline}\label{anexo-f.2.-desglose-tuxe9cnico-de-las-etapas-del-pipeline}

El ciclo de vida del pipeline se divide en dos trabajos secuenciales y
dependientes, los cuales ejecutan tareas críticas de aprovisionamiento,
autenticación y despliegue:

\subsubsection{\texorpdfstring{ANEXO F.2.1. Trabajo de Construcción y
Publicación
(\texttt{create-docker-image})}{ANEXO F.2.1. Trabajo de Construcción y Publicación (create-docker-image)}}\label{anexo-f.2.1.-trabajo-de-construcciuxf3n-y-publicaciuxf3n-create-docker-image}

Es el núcleo técnico de la automatización y consta de los siguientes
pasos detallados:

\begin{enumerate}
\def\labelenumi{\arabic{enumi}.}
\item
  \textbf{Clonación del Repositorio y Gestión de Submódulos:} Se utiliza
  la acción oficial \texttt{actions/checkout@v2}. Debido a la
  arquitectura desacoplada del proyecto, es indispensable configurar el
  parámetro
  \texttt{submodules:\ \textquotesingle{}recursive\textquotesingle{}}.
  Esta directiva instruye al agente para que descargue e integre de
  manera automática el repositorio y código de \texttt{evmaudit} dentro
  de la estructura de directorios del pipeline.
\item
  \textbf{Autenticación en el Registro de Contenedores (GHCR):} Mediante
  la acción \texttt{docker/login-action@v2}, el pipeline establece una
  conexión segura con el registro oficial de GitHub (\texttt{ghcr.io}).
  El proceso se autentica de forma dinámica utilizando el actor del
  ciclo de vida (\texttt{\$\{\{\ github.actor\ \}\}}) y un token de
  acceso seguro almacenado de forma cifrada en los secretos del
  repositorio bajo la clave
  \texttt{\$\{\{\ secrets.IMAGES\_TFM\_UNIR\ \}\}}.
\item
  \textbf{Normalización del Espacio de Nombres y Doble Etiquetado
  (\emph{Multi-tagging}):} Las especificaciones de Docker y el estándar
  de la Open Container Initiative (OCI) prohíben estrictamente el uso de
  caracteres en mayúscula para los nombres de las imágenes de
  contenedores. Para solucionar esto, el pipeline ejecuta un script en
  Bash que convierte dinámicamente el nombre del repositorio a
  minúsculas utilizando el comando
  \texttt{tr\ \textquotesingle{}{[}:upper:{]}\textquotesingle{}\ \textquotesingle{}{[}:lower:{]}\textquotesingle{}}.
  Posteriormente, se implementa una estrategia de doble etiquetado para
  optimizar la trazabilidad y la inmutabilidad:

  \begin{itemize}
  \tightlist
  \item
    \textbf{Etiqueta por SHA (}\texttt{IMAGE\_SHA}\textbf{):} Vincula de
    forma unívoca el contenedor con el hash del \emph{commit} específico
    de Git que originó la compilación
    (\texttt{\$\{\{\ github.sha\ \}\}}). Esto permite realizar
    auditorías retrospectivas y despliegues deterministas en caso de
    fallos.
  \item
    \textbf{Etiqueta de Última Versión
    (}\texttt{IMAGE\_LATEST}\textbf{):} Sobrescribe el puntero
    \texttt{:latest} con la versión más reciente del software que haya
    superado la fase de construcción con éxito.
  \end{itemize}
\item
  \textbf{Construcción y \emph{Push} Multiplataforma:} Se invoca de
  manera directa el comando \texttt{docker\ build}, forzando la
  compilación bajo la arquitectura destino
  \texttt{-\/-platform\ linux/amd64} utilizando el \texttt{Dockerfile}
  del proyecto como plano de construcción. Una vez generadas las
  imágenes locales con sus respectivas etiquetas, se ejecutan las
  instrucciones de empuje (\emph{push}) hacia el registro seguro de
  GitHub, quedando el artefacto disponible para su consumo.
\end{enumerate}

\subsubsection{ANEXO F.2.2. Trabajo de Despliegue (deploy) y
Limitaciones del
Entorno}\label{anexo-f.2.2.-trabajo-de-despliegue-deploy-y-limitaciones-del-entorno}

Para garantizar una separación formal de conceptos en la arquitectura de
la integración, el pipeline implementa un segundo trabajo secuencial
denominado \texttt{deploy}. Utilizando la directiva
\texttt{needs:\ create-docker-image}, se establece una restricción de
dependencia estricta: esta etapa no puede inicializarse si el
empaquetado y la publicación previa de la imagen en el registro han
fallado.

En la fase actual del proyecto, \textbf{esta etapa automatizada no ha
sido implementada de forma activa y opera estrictamente como un punto de
anclaje (\emph{placeholder}) arquitectónico}. La justificación de esta
decisión de diseño radica en las limitaciones técnicas del entorno de
alojamiento seleccionado para las pruebas de concepto. La
infraestructura de EVMAudit se despliega externamente en la plataforma
PaaS \textbf{Railway} utilizando su modalidad de suscripción gratuita.
Este nivel de servicio impone restricciones en las interfaces de
programación (APIs) y en el uso de \emph{webhooks}, impidiendo la
ejecución de despliegues totalmente automatizados (\emph{Automated CD
Triggers}) desencadenados de forma directa mediante agentes de terceros
como GitHub Actions.

Por consiguiente, el flujo de trabajo en este punto se limita a
verificar la integridad de la secuencia de comandos en consola, quedando
el aprovisionamiento de la infraestructura supeditado a los mecanismos
nativos de la plataforma de destino. En el \textbf{siguiente apartado
(X.5. Despliegue de la Infraestructura)}, se expondrá de manera más
amplia la configuración, el aprovisionamiento y las características
operativas de dicho entorno en Railway.


\chapter{Railway}
\section{ANEXO G. DESPLIEGUE DE LA INFRAESTRUCTURA EN LA NUBE
(RAILWAY)}\label{anexo-g.-despliegue-de-la-infraestructura-en-la-nube-railway}

Para validar la operatividad de EVMAudit en un entorno accesible y
simular un escenario de producción real, se ha procedido al despliegue
de la arquitectura contenedorizada en la plataforma de Plataforma como
Servicio (PaaS) \textbf{Railway}. A continuación, se detallan las
especificaciones del entorno, las restricciones técnicas de hardware
identificadas y las optimizaciones de ingeniería aplicadas en el código
fuente para garantizar la estabilidad del sistema.

\subsection{ANEXO G.1. Aprovisionamiento y Configuración del Entorno
Cloud}\label{anexo-g.1.-aprovisionamiento-y-configuraciuxf3n-del-entorno-cloud}

El proceso de despliegue en la infraestructura de la nube se ha
estructurado bajo las siguientes directrices operativas:

\begin{itemize}
\item
  \textbf{Selección del Nivel de Servicio:} La instancia se ha
  instanciado haciendo uso del nivel gratuito (\emph{Free Tier}) de la
  plataforma, el cual provee un crédito base de \$5 USD o un límite
  temporal de 30 días de cómputo.
\item
  \textbf{Aislamiento Regional:} Con el objetivo de minimizar la
  latencia de red en las peticiones HTTP y optimizar la transferencia de
  datos, se ha seleccionado la región europea con nodo central en
  \textbf{Ámsterdam (EU)}.
\item
  \textbf{Mapeo de Infraestructura y Orquestación:} El aprovisionamiento
  se realiza directamente vinculando el contenedor web a la imagen
  Docker compilada y almacenada en el registro de GitHub (ghcr.io),
  exponiendo de manera transparente la API del \emph{backend}
  desarrollada en FastAPI.
\item
  \textbf{Enrutamiento y Capa de Enlace (SSL):} La plataforma genera de
  manera dinámica un nombre de dominio completamente cualificado (FQDN)
  provisto de seguridad criptográfica TLS/SSL (HTTPS) para el acceso
  público a la interfaz de usuario:
  \url{https://evmaudit-production.up.railway.app/}.
\end{itemize}

\subsection{ANEXO G.2. Limitaciones de Hardware y el Problema del
Desbordamiento de Memoria
(OOM)}\label{anexo-g.2.-limitaciones-de-hardware-y-el-problema-del-desbordamiento-de-memoria-oom}

La modalidad gratuita de la plataforma PaaS impone restricciones
estrictas sobre los recursos de hardware asignados a cada contenedor,
parametrizados de la siguiente forma:

\begin{itemize}
\tightlist
\item
  \textbf{Capacidad de Cómputo (CPU):} 2 vCPU virtuales compartidas.
\item
  \textbf{Memoria Volátil (RAM):} 1 GB con un límite estricto de cuota
  (\emph{Plan Limit}).
\end{itemize}

Bajo un escenario de despliegue convencional en servidores dedicados o
infraestructura local, el pipeline de análisis de EVMAudit se ejecuta
sin restricciones debido a la disponibilidad de memoria elástica. Sin
embargo, en el entorno de la nube restringido, el motor de
\emph{fuzzing} basado en propiedades \textbf{Echidna} presenta un
problema crítico de arquitectura.

Echidna, al estar desarrollado en Haskell, requiere de forma nativa una
reserva inicial y un espacio de intercambio que supera con creces el
gigabyte de memoria RAM para gestionar el mapa de cobertura del binario
y la generación de casos de prueba. Al alcanzar el umbral crítico de 1
GB asignado por Railway, el demonio del sistema operativo anfitrión
(\emph{Kernel Out-of-Memory Killer}) destruía de manera abrupta el
contenedor para salvaguardar la integridad del nodo, provocando la caída
del servicio web y reportando un error de tipo \emph{OOM}.

\subsection{ANEXO G.3. Optimización del Sistema en Tiempo de Ejecución
(RTS) de
Haskell}\label{anexo-g.3.-optimizaciuxf3n-del-sistema-en-tiempo-de-ejecuciuxf3n-rts-de-haskell}

Para mitigar el desbordamiento de memoria sin alterar las capacidades
analíticas esenciales de la herramienta, se realizó una intervención a
nivel de código en el módulo de control del \emph{pipeline}
(run\_echidna). La solución consistió en inyectar directivas específicas
orientadas a reconfigurar los parámetros del \textbf{Runtime System
(RTS)} del compilador de Glasgow Haskell (GHC) empaquetados dentro del
binario de Echidna.

La estructura de invocación del proceso fue modificada e implementada en
Python mediante el siguiente diseño de argumentos:

\begin{Shaded}
\begin{Highlighting}[]
\NormalTok{command }\OperatorTok{=}\NormalTok{ [}
    \StringTok{"echidna"}\NormalTok{,}
\NormalTok{    contract\_path,}
    \StringTok{"{-}{-}contract"}\NormalTok{, contract\_name,}
    \StringTok{"{-}{-}format"}\NormalTok{, }\StringTok{"json"}\NormalTok{,}
    
    \CommentTok{\# Reducción del espacio de búsqueda del Fuzzer}
    \StringTok{"{-}{-}test{-}limit"}\NormalTok{, }\StringTok{"100"}\NormalTok{,  }
    
    \CommentTok{\# Optimización de la compilación inicial}
    \StringTok{"{-}{-}solc{-}args"}\NormalTok{, }\StringTok{"{-}{-}optimize{-}runs 0"}\NormalTok{, }
    
    \CommentTok{\# Apertura de las opciones del Runtime System de Haskell}
    \StringTok{"+RTS"}\NormalTok{,      }
    
    \CommentTok{\# Parámetros estrictos de control de memoria y concurrencia}
    \StringTok{"{-}M950m"}\NormalTok{,    }
    \StringTok{"{-}c"}\NormalTok{,        }
    \StringTok{"{-}N1"}\NormalTok{,       }
    
    \CommentTok{\# Cierre de las opciones RTS}
    \StringTok{"{-}RTS"}       
\NormalTok{]}
\end{Highlighting}
\end{Shaded}

A continuación se expone la justificación técnica detrás de los
modificadores inyectados:

\begin{itemize}
\tightlist
\item
  \textbf{Restricción de Ensayos (-\/-test-limit 100):} Al parametrizar
  el límite de pruebas en 100 iteraciones, se acota la profundidad del
  grafo de ejecución generado por el fuzzer. Esto reduce el consumo de
  memoria acumulativo a lo largo del tiempo de computación.
\item
  \textbf{Techo Infranqueable de Memoria (-M950m):} Esta directiva
  establece que el asignador de memoria de Haskell tiene prohibido
  estrictamente reclamar más de 950 megabytes del espacio de usuario. Al
  situar este límite ligeramente por debajo del gigabyte real de
  Railway, se evita que el sistema operativo de la nube elimine el
  proceso por exceder la cuota física.
\item
  \textbf{Recolección de Basura Agresiva (-c):} Activa un algoritmo de
  recolección de residuos más severo en el recolector de basura de
  Haskell (\emph{Garbage Collector}). En lugar de acumular objetos
  intermedios en la memoria RAM, el sistema libera los recursos
  obsoletos inmediatamente después de cada evaluación de propiedad.
\item
  \textbf{Concurrencia Confinada (-N1):} Limita la ejecución del entorno
  de ejecución a un único hilo de procesamiento, evitando la duplicación
  de estructuras de datos en memoria asociadas al paralelismo de hilos
  nativos.
\end{itemize}

Es importante destacar que la optimización descrita no se integró d
forma estática en la construcción del archivo Dockerfile (lo que habrí
alterado el comportamiento de la imagen base de manera permanente), sin
que se aplicó directamente sobre el código fuente desplegado en el
contenedor en ejecución (\emph{runtime}). Bajo condiciones de despliegue
convencionales en infraestructuras con escalabilidad elástica o recurso
dedicados de hardware, esta intervención técnica resultaría
completamente innecesaria, ya que la aplicación contaría con la memoria
suficiente para procesar el \emph{pipeline} por defecto. Por
consiguiente, esta modificación responde de manera estricta a un
mecanismo de mitigación ad hoc, implementado exclusivamente para sortear
las limitaciones físicas del entorno gratuito y evitar la interrupción
forzada del servicio web por falta de memoria.

\textbf{Conclusión del Despliegue:} La implementación de estas
salvaguardas de bajo nivel ha permitido que la aplicación web de
\textbf{EVMAudit} opere de manera completamente estable y fluida en la
nube. Pese a las severas restricciones del entorno gratuito de Railway,
el sistema es capaz de completar con éxito el pipeline completo de
auditoría en siete pasos sin registrar caídas en el servicio ni
excepciones por falta de recursos.

\end{document}